% =====================================================================
% appendix_target_generator.tex
% IEEE Access #Access-2026-29524 -- resubmission, Reviewer 1, point 2.
% Appendix documenting the deterministic target generator that the GNN
% is now positioned as EMULATING.
%
% Usage: % =====================================================================
% appendix_target_generator.tex
% IEEE Access #Access-2026-29524 -- resubmission, Reviewer 1, point 2.
% Appendix documenting the deterministic target generator that the GNN
% is now positioned as EMULATING.
%
% Usage: % =====================================================================
% appendix_target_generator.tex
% IEEE Access #Access-2026-29524 -- resubmission, Reviewer 1, point 2.
% Appendix documenting the deterministic target generator that the GNN
% is now positioned as EMULATING.
%
% Usage: % =====================================================================
% appendix_target_generator.tex
% IEEE Access #Access-2026-29524 -- resubmission, Reviewer 1, point 2.
% Appendix documenting the deterministic target generator that the GNN
% is now positioned as EMULATING.
%
% Usage: \input{REVISAO_R1/appendix_target_generator} immediately before
% the bibliography in main.tex. Requires amsmath and booktabs, both
% already loaded in the preamble (main.tex:11,16).
%
% PROVENANCE DISCIPLINE: every numerical value printed by this file is
% listed in the commented provenance table at the end, with the artefact
% field and the command that produced it. No number was read off a log.
% Artefact: REVISAO_R1/audit_target_generator.json
% Script:   REVISAO_R1/audit_target_generator.py
% =====================================================================

\section{Target Generation Model}
\label{app:target-generator}

\subsection{Software identification and scope}

The reference fields used throughout this study are not field measurements
but the output of a deterministic propagation generator written for this
project, and the contribution evaluated in this paper is
therefore the emulation of that generator at graph scale rather than the
prediction of measured radio conditions. No commercial or third-party
planning tool participates in the pipeline: there is no Atoll, ICS Telecom,
WinProp, TIREM or Longley--Rice component anywhere in the chain, and the
entire target field is produced by three Python modules and a post-processing corrector that depend, for
the propagation arithmetic itself, only on NumPy. Because those modules
carry no version string, they are identified here by the SHA-256 digest of
their contents, which is the only unambiguous version key available:
\texttt{generate\_\allowbreak{}realistic\_\allowbreak{}coverage\allowbreak{}.py} (\texttt{95ea0423\,2280\,6fd6},
13\,631~bytes) holds the propagation equations,
\texttt{01\_\allowbreak{}data/\allowbreak{}prepare\_\allowbreak{}transfer\_\allowbreak{}dataset\_\allowbreak{}v19.py}
(\texttt{45c16d43\,bc13\,a4bc}, 22\,917~bytes) applies them over the terrain
graph and writes the five supervised channels, and
\texttt{data\_\allowbreak{}raw/\allowbreak{}enrich\_\allowbreak{}rf\_\allowbreak{}targets\allowbreak{}.py} (\texttt{5d38012e\,c66e\,842d},
26\,012~bytes) replaces the two auxiliary channels with terrain-derived
attenuation terms. The corrector \texttt{contrafactual\_\allowbreak{}alvo\_\allowbreak{}completo\allowbreak{}.py}
(\texttt{ebeaf759\,2511\,fae3}, 10\,076~bytes) then rewrote the enriched files as
described in Section~\ref{app:corrector}; every result of the spatially blocked
evaluation uses the corrected files, whereas the comparisons quoted from the
earlier campaign use the enriched files before correction. Antenna
parameters come from a single ANATEL licensing extract,
\texttt{antenas\_\allowbreak{}interior\_\allowbreak{}sp\_\allowbreak{}final\allowbreak{}.csv}
(\texttt{ad9addb0\,c515\,a0b4}, 17\,763\,282~bytes). Everything reported in
this appendix is transcribed from those five files, and the equations below
are annotated with the source line at which each expression is implemented
so that the correspondence can be checked line by line.

\subsection{Base path loss}

The base term is a frequency-switched pair of empirical median-path-loss
models, both evaluated in decibels with distance in kilometres and
frequency in megahertz. The mobile-height correction factor for a
medium-to-small city is applied identically in both branches,
\begin{equation}
a(h_{m}) = \bigl(1.1\log_{10} f - 0.7\bigr)h_{m}
         - \bigl(1.56\log_{10} f - 0.8\bigr) ,
\label{eq:app-ahm}
\end{equation}
% generate_realistic_coverage.py:38 (Okumura branch) and :62 (COST-231 branch)
with the mobile height fixed at $h_{m}=1.5$~m for every node in the domain
(\texttt{prepare\_\allowbreak{}transfer\_\allowbreak{}dataset\_\allowbreak{}v19.py:\allowbreak{}210}). The Okumura--Hata urban
reference curve is
\begin{equation}
\begin{split}
L_{\mathrm{urb}}^{\mathrm{OH}} = {}& 69.55 + 26.16\log_{10} f
   - 13.82\log_{10} h_{b} - a(h_{m}) \\
 & + \bigl(44.9 - 6.55\log_{10} h_{b}\bigr)\log_{10} d ,
\end{split}
\label{eq:app-hata-urban}
\end{equation}
% generate_realistic_coverage.py:41-42
and the environment corrections applied to it are
\begin{equation}
L_{\mathrm{base}}^{\mathrm{OH}} =
\begin{cases}
L_{\mathrm{urb}}^{\mathrm{OH}}, & \text{urban},\\[2pt]
L_{\mathrm{urb}}^{\mathrm{OH}} - 2\bigl(\log_{10}\tfrac{f}{28}\bigr)^{2} - 5.4,
  & \text{suburban},\\[2pt]
L_{\mathrm{urb}}^{\mathrm{OH}} - 4.78\bigl(\log_{10} f\bigr)^{2} & \\
  \qquad {} + 18.33\log_{10} f - 40.94, & \text{rural}.
\end{cases}
\label{eq:app-hata-env}
\end{equation}

% generate_realistic_coverage.py:44-51
Above 1500~MHz the generator switches to the COST\,231 Hata extension,
\begin{equation}
\begin{split}
L_{\mathrm{base}}^{\mathrm{C231}} = {}& 46.3 + 33.9\log_{10} f
   - 13.82\log_{10} h_{b} - a(h_{m}) \\
 & + \bigl(44.9 - 6.55\log_{10} h_{b}\bigr)\log_{10} d + C_{m} ,
\end{split}
\label{eq:app-cost231}
\end{equation}
% generate_realistic_coverage.py:67-68; C_m at :65
where $C_{m}=3$~dB for the urban class and $C_{m}=0$~dB otherwise, and the
selection between the two branches is the hard threshold
$L_{\mathrm{base}} = L_{\mathrm{base}}^{\mathrm{OH}}$ if $f\le 1500$~MHz and
$L_{\mathrm{base}} = L_{\mathrm{base}}^{\mathrm{C231}}$ otherwise
(\texttt{generate\_\allowbreak{}realistic\_\allowbreak{}coverage\allowbreak{}.py:\allowbreak{}85--88}).

Three numerical guards inside these branches change the effective inputs
and must be declared because they are not part of either published model.
Frequency is clipped to $[150, 2000]$~MHz in the Okumura branch and to
$[1500, 2000]$~MHz in the COST\,231 branch
(\texttt{generate\_\allowbreak{}realistic\_\allowbreak{}coverage\allowbreak{}.py:\allowbreak{}34,58}), so that every carrier
above 2000~MHz enters the base term as if it were at exactly 2000~MHz,
while the vegetation term of Eq.~\eqref{eq:app-aveg} continues to receive
the true carrier (\texttt{generate\_\allowbreak{}realistic\_\allowbreak{}coverage\allowbreak{}.py:\allowbreak{}101}) and the two
components of the same sum are consequently evaluated at different
frequencies; in the four cities this clip is active for 27.4\,\% of the
Lins transmitters, 26.8\,\% of Campinas, 24.5\,\% of Sorocaba and 30.1\,\%
of Bauru, which means that for roughly a quarter of the field the base
frequency has been silently substituted. Distance is floored at $10^{-3}$~km,
that is at one metre (\texttt{generate\_\allowbreak{}realistic\_\allowbreak{}coverage\allowbreak{}.py:\allowbreak{}35,59}),
three decades below the 1~km lower bound of the Okumura--Hata validity
domain. Base-station height enters unclipped, and 5.1\,\% of the Lins
transmitters, 23.3\,\% of Campinas, 18.9\,\% of Sorocaba and 6.7\,\% of
Bauru sit below the 30~m floor of that same domain. The generator is
therefore an extrapolation of both empirical models over part of its input
range, and the reported accuracy figures should be read as fidelity to this
particular extrapolation rather than as agreement with the published
models inside their stated envelopes.

\subsection{Clutter and environment class}

What the generator calls the propagation environment is not a clutter
parameter in the usual sense, and this is the single point in the model
most likely to be misread. There is no land-cover raster, no morphology
database, no building or clutter-height layer, and no per-pixel clutter
correction anywhere in the chain; the class is assigned once per
transmitter by a threshold on its own carrier frequency, urban when
$f>1800$~MHz, suburban when $800<f\le 1800$~MHz and rural when
$f\le 800$~MHz (\texttt{prepare\_\allowbreak{}transfer\_\allowbreak{}dataset\_\allowbreak{}v19.py:\allowbreak{}84--86}). The
consequence is that the environment label is a deterministic function of
frequency and carries no independent geographic information, so that
Eq.~\eqref{eq:app-hata-env} does not describe the ground under a given node
but only the band of the transmitter serving it. A verifiable corollary of
this coupling is that the urban class can never reach the Okumura branch
and the rural class can never reach COST\,231, which the transmitter records
confirm exactly: zero transmitters fall in either impossible combination in any of
the four cities. Terrain clutter is therefore unmodelled, and a frequency-indexed
correction constant stands in its place.

\subsection{Vegetation attenuation}

Vegetation enters through a single additive term parameterised by NDVI
taken raw from column 12 of the terrain feature matrix
(\texttt{prepare\_\allowbreak{}transfer\_\allowbreak{}dataset\_\allowbreak{}v19.py:\allowbreak{}140}), which the code maps to a
density and then to an effective path depth,
\begin{equation}
\rho_{v} = \operatorname{clip}\!\left(\frac{\mathrm{NDVI}+1}{2},\,0,\,1\right),
\qquad
d_{v} = 0.3\, \rho_{v}\, d_{m} ,
\label{eq:app-vegdepth}
\end{equation}
% generate_realistic_coverage.py:95,98 -- d_m in metres
so that at most 30\,\% of the geometric path is declared to lie inside
vegetation regardless of where the vegetation actually is. The attenuation
follows the ITU-R P.833 functional form for a single vegetation entry,
\begin{equation}
A_{\mathrm{veg}} = \operatorname{clip}\!\bigl(
0.25 f^{0.39} \min(d_{v}, 400)^{0.25},\, 0,\, 40 \bigr)\ [\mathrm{dB}] ,
\label{eq:app-aveg}
\end{equation}
% generate_realistic_coverage.py:101-102 -- f in MHz, d_v in metres
with frequency in megahertz and depth in metres, a depth cap of 400~m and
an attenuation cap of 40~dB. The depth cap is the dominant feature of this
term in practice: for a node with $\mathrm{NDVI}=0.5$ the cap is reached at
a geometric distance of about 1778~m, beyond which $A_{\mathrm{veg}}$ is
constant in distance and depends only on frequency, taking the values
15.87~dB at 900~MHz, 20.80~dB at 1800~MHz and 24.00~dB at 2600~MHz.
Because the argument of the quarter-power is a product of distance and
NDVI, the term is also strongly distance-driven below the cap, reaching
7.73~dB at 100~m and 13.74~dB at 1~km for the same NDVI at 900~MHz, so that
most of its dynamic range is geometry rather than vegetation. The 40~dB
clip is unreachable in this dataset: at the highest carrier present in the
four cities, 3550~MHz, and with the depth cap and NDVI both saturated, the
term attains only 27.10~dB.

\subsection{Terrain attenuation}

The terrain term is the component of the generator that most needs an
explicit and uncomfortable statement, because its name does not describe
its content. Formally it is a piecewise-linear function of a dimensionless
slope variable,
\begin{equation}
A_{\mathrm{ter}} = \operatorname{clip}\!\Bigl(
20\,s + 30\,\max(s-0.1,\,0),\; 0,\; 30 \Bigr)\ [\mathrm{dB}] ,
\label{eq:app-aterrain}
\end{equation}
% generate_realistic_coverage.py:109-116
with $s$ the local slope expressed as a rise-over-run ratio. The V19 stage
of the pipeline did not read $s$ from the digital elevation model: it drew
it independently for each node from a uniform distribution on $[0.02, 0.15]$
by an unseeded call to the NumPy generator
(\texttt{prepare\_\allowbreak{}transfer\_\allowbreak{}dataset\_\allowbreak{}v19.py:\allowbreak{}141}, and identically at
\texttt{generate\_\allowbreak{}realistic\_\allowbreak{}coverage\allowbreak{}.py:\allowbreak{}167}), which made
Eq.~\eqref{eq:app-aterrain} an independent random offset per node between
0.4~dB and 4.5~dB, with mean 1.99~dB and standard deviation 1.15~dB over the
support of that distribution. The corrector of Section~\ref{app:corrector}
replaces the draw with the node's DEM slope, stored as feature~1 of the
released files, and recomputes the total-loss and RSSI channels by exact
difference, so that in every reference field used for the blocked evaluation
Eq.~\eqref{eq:app-aterrain} is a deterministic function of terrain.
Recomputing the corrected total-loss channel from the features stored in the
released files reproduces it with a mean absolute error of at most
$3.02 \times 10^{-4}$~dB per cell over 150\,000 covered nodes in five cells.

\subsection{Total loss, link budget and serving cell}

The three components are summed without interaction terms,
\begin{equation}
\begin{split}
L_{\mathrm{tot}}(i,j) = {}& L_{\mathrm{base}}(d_{ij}, f_{j}, h_{b,j}, e_{j}) \\
 & + A_{\mathrm{veg}}(d_{ij}, \mathrm{NDVI}_{i})
   + A_{\mathrm{ter}}(s_{i}) ,
\end{split}
\label{eq:app-ltot}
\end{equation}
% generate_realistic_coverage.py:121
and converted to a received level by a two-term link budget that contains
no cable, connector, body or feeder loss and no receive antenna gain,
\begin{equation}
\mathrm{RSSI}(i,j) = P_{j} + G_{j} - L_{\mathrm{tot}}(i,j)
\quad [\mathrm{dBm}] ,
\label{eq:app-rssi}
\end{equation}
% prepare_transfer_dataset_v19.py:215
where $P_{j}=10\log_{10}(1000\,P_{j}^{\mathrm{W}})$ is the transmitter power
in dBm (\texttt{prepare\_\allowbreak{}transfer\_\allowbreak{}dataset\_\allowbreak{}v19.py:\allowbreak{}81}) and $G_{j}$ is the
catalogue antenna gain in dBi, applied as a scalar in every direction. The
serving cell is then the strongest received level among the transmitters
whose planar distance to the node is within the 30~km inclusion radius,
\begin{equation}
s(i) = \!\!\operatorname*{arg\,max}_{j \,:\, d_{ij} \le 30\,\mathrm{km}}\!\!
\mathrm{RSSI}(i,j),
\quad
L_{i} = L_{\mathrm{tot}}\bigl(i, s(i)\bigr) ,
\label{eq:app-serving}
\end{equation}
% prepare_transfer_dataset_v19.py:220-228; radius at :46
which is a pure best-server rule with no handover hysteresis, no load
balancing, no cell-selection offset and no minimum-received-level
qualification. Distances are planar rather than geodesic, computed as
$\Delta y = 111000\,\Delta\varphi$ and
$\Delta x = 111000\,\Delta\lambda\cos\varphi$ in metres
(\texttt{prepare\_\allowbreak{}transfer\_\allowbreak{}dataset\_\allowbreak{}v19.py:\allowbreak{}183--185}); the constant
111000~m per degree understates the WGS-84 mean meridian degree of
111194.93~m by 0.175\,\%, a bias that is negligible against the decibel
scale of the model but is nonetheless a systematic contraction of every
distance in the dataset. Nodes with no transmitter inside the radius retain
the sentinel initialisation of $-200$~dBm and 300~dB, and any node whose
best level falls below $-150$~dBm is subsequently overwritten with a
constant $-110$~dBm (\texttt{prepare\_\allowbreak{}transfer\_\allowbreak{}dataset\_\allowbreak{}v19.py:\allowbreak{}274}),
which introduces a point mass in the RSSI distribution that must be
excluded before any distributional claim about the target field is made.

\subsection{Coverage definition}

Coverage is a hard indicator rather than a probability, defined on the
serving-cell received level by a single threshold,
\begin{equation}
c_{i} = \mathbf{1}\bigl[\mathrm{RSSI}_{i} > -85\ \mathrm{dBm}\bigr] ,
\label{eq:app-coverage}
\end{equation}
% prepare_transfer_dataset_v19.py:275
so that $H(c_{i}\mid \mathrm{RSSI}_{i}) = 0$ exactly and the fifth
supervised channel carries no information beyond the fourth. There is no
fading margin, no outage probability, no shadowing standard deviation and
no receiver sensitivity or noise-figure calculation entering this
definition; the threshold is a fixed design constant of the generator. The
learned coverage channel is accordingly a derived classification of the
received level rather than an independent coverage probability.

\subsection{Antenna configuration}

Table~\ref{tab:app-antenna-params} separates the antenna quantities the
generator actually consumes from those that exist in the licensing extract
and are discarded, and Table~\ref{tab:app-antenna-inventory} gives the
resulting per-city inventory. The essential point for a reader assessing
the physical realism of the target is that the radiation pattern is
isotropic: azimuth, mechanical and electrical downtilt, half-power
beamwidth, front-to-back ratio and polarisation are all present in the
source records and none is read by the loader
(\texttt{prepare\_\allowbreak{}transfer\_\allowbreak{}dataset\_\allowbreak{}v19.py:\allowbreak{}73--79}). Because sectors of the
same site share coordinates, discarding azimuth turns each site into a
stack of co-located omnidirectional sources that differ only in frequency,
power, gain and height, with a mean of 13.5 and a maximum of 21 such
sources per site in Lins and a maximum of 49 in Campinas; the best-server
rule of Eq.~\eqref{eq:app-serving} then reduces each stack to whichever
member happens to maximise $P_{j}+G_{j}-L_{\mathrm{tot}}$. The loader also
retains catalogue gains of exactly 0~dBi where the licensing record carries
no value, affecting 32 transmitters in Campinas, 4 in Sorocaba and 1 in
Bauru, and it does not apply the 12~dBi substitution that the standalone
coverage script uses for the same situation
(\texttt{generate\_\allowbreak{}realistic\_\allowbreak{}coverage\allowbreak{}.py:\allowbreak{}191}), so the two entry points to
the same physics are not numerically identical.

\begin{table}[!t]
\caption{Antenna quantities in the target generator, grouped by whether
the model consumes them. Column names in typewriter font are fields of the
ANATEL licensing extract. The second group is the reason the radiation
pattern is isotropic: those fields are populated in the records and the
loader never reads them.}
\label{tab:app-antenna-params}
\centering
\footnotesize
\begin{tabular}{@{}ll@{}}
\toprule
Quantity & Source and unit in the model \\
\midrule
\multicolumn{2}{@{}l}{\textit{Used}} \\
Carrier frequency $f$   & \texttt{FreqTxMHz}, MHz \\
Transmit power $P$      & $10\log_{10}(1000P^{\mathrm{W}})$, dBm \\
Antenna gain $G$        & \texttt{GanhoAntena}, dBi, isotropic \\
Antenna height $h_{b}$  & \texttt{AlturaAntena}, m AGL \\
Site coordinates        & \texttt{Latitude}, \texttt{Longitude}, deg \\
Environment class $e$   & derived from $f$ alone \\
Mobile height $h_{m}$   & fixed at 1.5\,m for all nodes \\
\midrule
\multicolumn{2}{@{}l}{\textit{Present in the records but never read}} \\
Azimuth                 & \texttt{Azimute}, deg \\
Downtilt                & \texttt{AnguloElevacao}, deg \\
Half-power beamwidth    & \texttt{AnguloMeiaPotenciaAntena}, deg \\
Front-to-back ratio     & \texttt{FrenteCostaAntena}, dB \\
Polarisation            & \texttt{Polarizacao} \\
\midrule
\multicolumn{2}{@{}l}{\textit{Absent from the model altogether}} \\
Feeder, connector, body loss & --- \\
Receive antenna gain         & --- \\
Noise figure, thermal noise  & --- \\
Interference, SINR           & --- \\
Fading, shadowing margin     & --- \\
\bottomrule
\end{tabular}
\end{table}

\begin{table}[!t]
\caption{Per-city transmitter inventory after the generator's validity
filter ($f>0$ and $P^{\mathrm{W}}>0$). Sites are unique coordinate pairs.
The last column is the share of transmitters whose carrier exceeds the
2000\,MHz internal clip.}
\label{tab:app-antenna-inventory}
\centering
\footnotesize
\setlength{\tabcolsep}{4pt}
\begin{tabular}{@{}lrrccrr@{}}
\toprule
City & Tx & Sites & $f$ (MHz) & $P$ (dBm) & $h_{b}$ (m) & Clip \\
\midrule
Lins     &  175 &  13 & 788--2680 & 41.6--49.0 & 21--75 & 27.4\,\% \\
Campinas & 4252 & 350 & 778--3550 & 33.0--53.0 &  3--74 & 26.8\,\% \\
Sorocaba & 1398 & 120 & 788--2680 & 35.1--49.3 &  3--76 & 24.5\,\% \\
Bauru    &  794 &  75 & 788--2680 & 40.0--49.0 &  3--90 & 30.1\,\% \\
\bottomrule
\end{tabular}
\end{table}

\subsection{The two auxiliary channels}

The five supervised channels written by the V19 pipeline are total path
loss, two placeholder channels, serving-cell RSSI and the coverage
indicator; the placeholders are written as zeros at that stage
(\texttt{prepare\_\allowbreak{}transfer\_\allowbreak{}dataset\_\allowbreak{}v19.py:\allowbreak{}276--282}) and are then filled
by the enrichment step, which is part of the pipeline for every dataset
used in this study. That step computes a vegetation shadow margin
following ITU-R P.833 from the Sentinel-2 NDVI of the node and a
knife-edge diffraction loss following ITU-R P.526 from the DEM-derived
terrain ruggedness index and roughness
(\texttt{data\_\allowbreak{}raw/\allowbreak{}enrich\_\allowbreak{}rf\_\allowbreak{}targets\allowbreak{}.py:\allowbreak{}336--396,533--534}), using a
specific attenuation $\gamma = 0.2\,\mathrm{NDVI}\,(f_{\mathrm{GHz}}/1.8)^{0.4}$
in dB/m over a depth $d_{v} = d\,\operatorname{clip}(\mathrm{NDVI}-0.5,0,1)$
capped at 40~dB, and a single knife-edge obstacle placed at the path
midpoint with an obstruction height
$h_{\mathrm{obs}} = \operatorname{clip}(0.5\,\mathrm{TRI}_{\mathrm{m}} + 0.5\,r_{\mathrm{m}},\,0,\,80)$
in metres, also capped at 40~dB, with roughness $r_{\mathrm{m}}$ in metres after
correction. These two channels are therefore the point
at which the digital elevation model and the vegetation index enter the
reference field as physically labelled attenuation terms, and they are the
targets against which the decoder's vegetation and terrain heads are
trained and scored. Two properties of the enrichment must be stated because
they differ from the geometry of the total-loss channel: the distance it
uses is the distance to the nearest transmitter rather than to the serving
cell of Eq.~\eqref{eq:app-serving}
(\texttt{enrich\_\allowbreak{}rf\_\allowbreak{}targets\allowbreak{}.py:\allowbreak{}546}), and the frequency is a single
global 1800~MHz (\texttt{FREQ\_\allowbreak{}MHZ\_\allowbreak{}DEFAULT}, line 38) rather than the
per-transmitter carrier. The auxiliary channels are consequently
consistent with each other but evaluated under a different link geometry
and frequency from the total-loss channel, and the manuscript states this
explicitly where the three channels are compared.

\subsection{Target corrector}
\label{app:corrector}

Three defects of the enriched files were corrected before the spatially
blocked evaluation, all by \texttt{contrafactual\_\allowbreak{}alvo\_\allowbreak{}completo\allowbreak{}.py}, which
writes new files and leaves the originals untouched. First, the terrain term
of the total-loss channel used the random slope draw described above; the
corrector substitutes the node's DEM slope and updates the total-loss and
RSSI channels by exact difference, which is possible because
Eq.~\eqref{eq:app-aterrain} depends only on the node and does not change which
transmitter serves it. Second, the enrichment step read the topographic
position index column where the diffraction term expects the terrain
ruggedness index; the corrector reads the ruggedness column. Third, the
obstruction height added ruggedness in metres to roughness in normalised
units, whose scale differs by city; the corrector converts roughness to
metres with the per-city elevation standard deviation. The diffraction
channel is then recomputed with the original enrichment function. Sentinel
nodes keep their sentinel values in the total-loss and RSSI channels. The
16 corrected reference fields are released with their SHA-256 digests.

\subsection{What the generator does not model}

An emulation claim is only as strong as the honesty of the description of
what is being emulated, and several mechanisms that a radio-planning reader
would expect are simply absent here. There is no ray tracing and no terrain
profile along the link: the base term never sees an elevation sample
between transmitter and receiver, so diffraction over real obstacles,
Fresnel-zone clearance, line-of-sight testing, multipath and ground
reflection have no representation, and the knife-edge term of the
enrichment substitutes a single midpoint obstacle whose height is a
ruggedness proxy rather than a profile maximum. There is no clutter
database, as discussed above, and no building layer, so the urban,
suburban and rural corrections act as three global constants selected by
carrier frequency. There is no stochastic component of any kind in the
delivered path loss: the docstring of the propagation routine announces
log-normal fading (\texttt{generate\_\allowbreak{}realistic\_\allowbreak{}coverage\allowbreak{}.py:\allowbreak{}80}) but no
random term appears in the body of the function, and the only randomness in
the pipeline was the terrain slope draw described above, which the corrector
removes from the released fields. There is no noise floor, no
receiver noise figure, no interference from neighbouring cells and
therefore no SINR, so the received level of Eq.~\eqref{eq:app-rssi} is a
pure two-term link budget and coverage is a threshold on that budget rather
than a link-quality decision. Finally, the geometry is two-dimensional and
planar: node and transmitter altitudes never enter, the receiver is fixed
at 1.5~m above an unspecified reference surface, Earth curvature is
ignored, and the distance uses a constant metres-per-degree factor. Stating
these absences precisely is what allows the reformulated contribution to be
defended, because the quantity the network reproduces is a well-defined and
fully specified analytical field, and the accuracy figures in the body of
the paper measure emulation fidelity to that field and nothing more.
 immediately before
% the bibliography in main.tex. Requires amsmath and booktabs, both
% already loaded in the preamble (main.tex:11,16).
%
% PROVENANCE DISCIPLINE: every numerical value printed by this file is
% listed in the commented provenance table at the end, with the artefact
% field and the command that produced it. No number was read off a log.
% Artefact: REVISAO_R1/audit_target_generator.json
% Script:   REVISAO_R1/audit_target_generator.py
% =====================================================================

\section{Target Generation Model}
\label{app:target-generator}

\subsection{Software identification and scope}

The reference fields used throughout this study are not field measurements
but the output of a deterministic propagation generator written for this
project, and the contribution evaluated in this paper is
therefore the emulation of that generator at graph scale rather than the
prediction of measured radio conditions. No commercial or third-party
planning tool participates in the pipeline: there is no Atoll, ICS Telecom,
WinProp, TIREM or Longley--Rice component anywhere in the chain, and the
entire target field is produced by three Python modules and a post-processing corrector that depend, for
the propagation arithmetic itself, only on NumPy. Because those modules
carry no version string, they are identified here by the SHA-256 digest of
their contents, which is the only unambiguous version key available:
\texttt{generate\_\allowbreak{}realistic\_\allowbreak{}coverage\allowbreak{}.py} (\texttt{95ea0423\,2280\,6fd6},
13\,631~bytes) holds the propagation equations,
\texttt{01\_\allowbreak{}data/\allowbreak{}prepare\_\allowbreak{}transfer\_\allowbreak{}dataset\_\allowbreak{}v19.py}
(\texttt{45c16d43\,bc13\,a4bc}, 22\,917~bytes) applies them over the terrain
graph and writes the five supervised channels, and
\texttt{data\_\allowbreak{}raw/\allowbreak{}enrich\_\allowbreak{}rf\_\allowbreak{}targets\allowbreak{}.py} (\texttt{5d38012e\,c66e\,842d},
26\,012~bytes) replaces the two auxiliary channels with terrain-derived
attenuation terms. The corrector \texttt{contrafactual\_\allowbreak{}alvo\_\allowbreak{}completo\allowbreak{}.py}
(\texttt{ebeaf759\,2511\,fae3}, 10\,076~bytes) then rewrote the enriched files as
described in Section~\ref{app:corrector}; every result of the spatially blocked
evaluation uses the corrected files, whereas the comparisons quoted from the
earlier campaign use the enriched files before correction. Antenna
parameters come from a single ANATEL licensing extract,
\texttt{antenas\_\allowbreak{}interior\_\allowbreak{}sp\_\allowbreak{}final\allowbreak{}.csv}
(\texttt{ad9addb0\,c515\,a0b4}, 17\,763\,282~bytes). Everything reported in
this appendix is transcribed from those five files, and the equations below
are annotated with the source line at which each expression is implemented
so that the correspondence can be checked line by line.

\subsection{Base path loss}

The base term is a frequency-switched pair of empirical median-path-loss
models, both evaluated in decibels with distance in kilometres and
frequency in megahertz. The mobile-height correction factor for a
medium-to-small city is applied identically in both branches,
\begin{equation}
a(h_{m}) = \bigl(1.1\log_{10} f - 0.7\bigr)h_{m}
         - \bigl(1.56\log_{10} f - 0.8\bigr) ,
\label{eq:app-ahm}
\end{equation}
% generate_realistic_coverage.py:38 (Okumura branch) and :62 (COST-231 branch)
with the mobile height fixed at $h_{m}=1.5$~m for every node in the domain
(\texttt{prepare\_\allowbreak{}transfer\_\allowbreak{}dataset\_\allowbreak{}v19.py:\allowbreak{}210}). The Okumura--Hata urban
reference curve is
\begin{equation}
\begin{split}
L_{\mathrm{urb}}^{\mathrm{OH}} = {}& 69.55 + 26.16\log_{10} f
   - 13.82\log_{10} h_{b} - a(h_{m}) \\
 & + \bigl(44.9 - 6.55\log_{10} h_{b}\bigr)\log_{10} d ,
\end{split}
\label{eq:app-hata-urban}
\end{equation}
% generate_realistic_coverage.py:41-42
and the environment corrections applied to it are
\begin{equation}
L_{\mathrm{base}}^{\mathrm{OH}} =
\begin{cases}
L_{\mathrm{urb}}^{\mathrm{OH}}, & \text{urban},\\[2pt]
L_{\mathrm{urb}}^{\mathrm{OH}} - 2\bigl(\log_{10}\tfrac{f}{28}\bigr)^{2} - 5.4,
  & \text{suburban},\\[2pt]
L_{\mathrm{urb}}^{\mathrm{OH}} - 4.78\bigl(\log_{10} f\bigr)^{2} & \\
  \qquad {} + 18.33\log_{10} f - 40.94, & \text{rural}.
\end{cases}
\label{eq:app-hata-env}
\end{equation}

% generate_realistic_coverage.py:44-51
Above 1500~MHz the generator switches to the COST\,231 Hata extension,
\begin{equation}
\begin{split}
L_{\mathrm{base}}^{\mathrm{C231}} = {}& 46.3 + 33.9\log_{10} f
   - 13.82\log_{10} h_{b} - a(h_{m}) \\
 & + \bigl(44.9 - 6.55\log_{10} h_{b}\bigr)\log_{10} d + C_{m} ,
\end{split}
\label{eq:app-cost231}
\end{equation}
% generate_realistic_coverage.py:67-68; C_m at :65
where $C_{m}=3$~dB for the urban class and $C_{m}=0$~dB otherwise, and the
selection between the two branches is the hard threshold
$L_{\mathrm{base}} = L_{\mathrm{base}}^{\mathrm{OH}}$ if $f\le 1500$~MHz and
$L_{\mathrm{base}} = L_{\mathrm{base}}^{\mathrm{C231}}$ otherwise
(\texttt{generate\_\allowbreak{}realistic\_\allowbreak{}coverage\allowbreak{}.py:\allowbreak{}85--88}).

Three numerical guards inside these branches change the effective inputs
and must be declared because they are not part of either published model.
Frequency is clipped to $[150, 2000]$~MHz in the Okumura branch and to
$[1500, 2000]$~MHz in the COST\,231 branch
(\texttt{generate\_\allowbreak{}realistic\_\allowbreak{}coverage\allowbreak{}.py:\allowbreak{}34,58}), so that every carrier
above 2000~MHz enters the base term as if it were at exactly 2000~MHz,
while the vegetation term of Eq.~\eqref{eq:app-aveg} continues to receive
the true carrier (\texttt{generate\_\allowbreak{}realistic\_\allowbreak{}coverage\allowbreak{}.py:\allowbreak{}101}) and the two
components of the same sum are consequently evaluated at different
frequencies; in the four cities this clip is active for 27.4\,\% of the
Lins transmitters, 26.8\,\% of Campinas, 24.5\,\% of Sorocaba and 30.1\,\%
of Bauru, which means that for roughly a quarter of the field the base
frequency has been silently substituted. Distance is floored at $10^{-3}$~km,
that is at one metre (\texttt{generate\_\allowbreak{}realistic\_\allowbreak{}coverage\allowbreak{}.py:\allowbreak{}35,59}),
three decades below the 1~km lower bound of the Okumura--Hata validity
domain. Base-station height enters unclipped, and 5.1\,\% of the Lins
transmitters, 23.3\,\% of Campinas, 18.9\,\% of Sorocaba and 6.7\,\% of
Bauru sit below the 30~m floor of that same domain. The generator is
therefore an extrapolation of both empirical models over part of its input
range, and the reported accuracy figures should be read as fidelity to this
particular extrapolation rather than as agreement with the published
models inside their stated envelopes.

\subsection{Clutter and environment class}

What the generator calls the propagation environment is not a clutter
parameter in the usual sense, and this is the single point in the model
most likely to be misread. There is no land-cover raster, no morphology
database, no building or clutter-height layer, and no per-pixel clutter
correction anywhere in the chain; the class is assigned once per
transmitter by a threshold on its own carrier frequency, urban when
$f>1800$~MHz, suburban when $800<f\le 1800$~MHz and rural when
$f\le 800$~MHz (\texttt{prepare\_\allowbreak{}transfer\_\allowbreak{}dataset\_\allowbreak{}v19.py:\allowbreak{}84--86}). The
consequence is that the environment label is a deterministic function of
frequency and carries no independent geographic information, so that
Eq.~\eqref{eq:app-hata-env} does not describe the ground under a given node
but only the band of the transmitter serving it. A verifiable corollary of
this coupling is that the urban class can never reach the Okumura branch
and the rural class can never reach COST\,231, which the transmitter records
confirm exactly: zero transmitters fall in either impossible combination in any of
the four cities. Terrain clutter is therefore unmodelled, and a frequency-indexed
correction constant stands in its place.

\subsection{Vegetation attenuation}

Vegetation enters through a single additive term parameterised by NDVI
taken raw from column 12 of the terrain feature matrix
(\texttt{prepare\_\allowbreak{}transfer\_\allowbreak{}dataset\_\allowbreak{}v19.py:\allowbreak{}140}), which the code maps to a
density and then to an effective path depth,
\begin{equation}
\rho_{v} = \operatorname{clip}\!\left(\frac{\mathrm{NDVI}+1}{2},\,0,\,1\right),
\qquad
d_{v} = 0.3\, \rho_{v}\, d_{m} ,
\label{eq:app-vegdepth}
\end{equation}
% generate_realistic_coverage.py:95,98 -- d_m in metres
so that at most 30\,\% of the geometric path is declared to lie inside
vegetation regardless of where the vegetation actually is. The attenuation
follows the ITU-R P.833 functional form for a single vegetation entry,
\begin{equation}
A_{\mathrm{veg}} = \operatorname{clip}\!\bigl(
0.25 f^{0.39} \min(d_{v}, 400)^{0.25},\, 0,\, 40 \bigr)\ [\mathrm{dB}] ,
\label{eq:app-aveg}
\end{equation}
% generate_realistic_coverage.py:101-102 -- f in MHz, d_v in metres
with frequency in megahertz and depth in metres, a depth cap of 400~m and
an attenuation cap of 40~dB. The depth cap is the dominant feature of this
term in practice: for a node with $\mathrm{NDVI}=0.5$ the cap is reached at
a geometric distance of about 1778~m, beyond which $A_{\mathrm{veg}}$ is
constant in distance and depends only on frequency, taking the values
15.87~dB at 900~MHz, 20.80~dB at 1800~MHz and 24.00~dB at 2600~MHz.
Because the argument of the quarter-power is a product of distance and
NDVI, the term is also strongly distance-driven below the cap, reaching
7.73~dB at 100~m and 13.74~dB at 1~km for the same NDVI at 900~MHz, so that
most of its dynamic range is geometry rather than vegetation. The 40~dB
clip is unreachable in this dataset: at the highest carrier present in the
four cities, 3550~MHz, and with the depth cap and NDVI both saturated, the
term attains only 27.10~dB.

\subsection{Terrain attenuation}

The terrain term is the component of the generator that most needs an
explicit and uncomfortable statement, because its name does not describe
its content. Formally it is a piecewise-linear function of a dimensionless
slope variable,
\begin{equation}
A_{\mathrm{ter}} = \operatorname{clip}\!\Bigl(
20\,s + 30\,\max(s-0.1,\,0),\; 0,\; 30 \Bigr)\ [\mathrm{dB}] ,
\label{eq:app-aterrain}
\end{equation}
% generate_realistic_coverage.py:109-116
with $s$ the local slope expressed as a rise-over-run ratio. The V19 stage
of the pipeline did not read $s$ from the digital elevation model: it drew
it independently for each node from a uniform distribution on $[0.02, 0.15]$
by an unseeded call to the NumPy generator
(\texttt{prepare\_\allowbreak{}transfer\_\allowbreak{}dataset\_\allowbreak{}v19.py:\allowbreak{}141}, and identically at
\texttt{generate\_\allowbreak{}realistic\_\allowbreak{}coverage\allowbreak{}.py:\allowbreak{}167}), which made
Eq.~\eqref{eq:app-aterrain} an independent random offset per node between
0.4~dB and 4.5~dB, with mean 1.99~dB and standard deviation 1.15~dB over the
support of that distribution. The corrector of Section~\ref{app:corrector}
replaces the draw with the node's DEM slope, stored as feature~1 of the
released files, and recomputes the total-loss and RSSI channels by exact
difference, so that in every reference field used for the blocked evaluation
Eq.~\eqref{eq:app-aterrain} is a deterministic function of terrain.
Recomputing the corrected total-loss channel from the features stored in the
released files reproduces it with a mean absolute error of at most
$3.02 \times 10^{-4}$~dB per cell over 150\,000 covered nodes in five cells.

\subsection{Total loss, link budget and serving cell}

The three components are summed without interaction terms,
\begin{equation}
\begin{split}
L_{\mathrm{tot}}(i,j) = {}& L_{\mathrm{base}}(d_{ij}, f_{j}, h_{b,j}, e_{j}) \\
 & + A_{\mathrm{veg}}(d_{ij}, \mathrm{NDVI}_{i})
   + A_{\mathrm{ter}}(s_{i}) ,
\end{split}
\label{eq:app-ltot}
\end{equation}
% generate_realistic_coverage.py:121
and converted to a received level by a two-term link budget that contains
no cable, connector, body or feeder loss and no receive antenna gain,
\begin{equation}
\mathrm{RSSI}(i,j) = P_{j} + G_{j} - L_{\mathrm{tot}}(i,j)
\quad [\mathrm{dBm}] ,
\label{eq:app-rssi}
\end{equation}
% prepare_transfer_dataset_v19.py:215
where $P_{j}=10\log_{10}(1000\,P_{j}^{\mathrm{W}})$ is the transmitter power
in dBm (\texttt{prepare\_\allowbreak{}transfer\_\allowbreak{}dataset\_\allowbreak{}v19.py:\allowbreak{}81}) and $G_{j}$ is the
catalogue antenna gain in dBi, applied as a scalar in every direction. The
serving cell is then the strongest received level among the transmitters
whose planar distance to the node is within the 30~km inclusion radius,
\begin{equation}
s(i) = \!\!\operatorname*{arg\,max}_{j \,:\, d_{ij} \le 30\,\mathrm{km}}\!\!
\mathrm{RSSI}(i,j),
\quad
L_{i} = L_{\mathrm{tot}}\bigl(i, s(i)\bigr) ,
\label{eq:app-serving}
\end{equation}
% prepare_transfer_dataset_v19.py:220-228; radius at :46
which is a pure best-server rule with no handover hysteresis, no load
balancing, no cell-selection offset and no minimum-received-level
qualification. Distances are planar rather than geodesic, computed as
$\Delta y = 111000\,\Delta\varphi$ and
$\Delta x = 111000\,\Delta\lambda\cos\varphi$ in metres
(\texttt{prepare\_\allowbreak{}transfer\_\allowbreak{}dataset\_\allowbreak{}v19.py:\allowbreak{}183--185}); the constant
111000~m per degree understates the WGS-84 mean meridian degree of
111194.93~m by 0.175\,\%, a bias that is negligible against the decibel
scale of the model but is nonetheless a systematic contraction of every
distance in the dataset. Nodes with no transmitter inside the radius retain
the sentinel initialisation of $-200$~dBm and 300~dB, and any node whose
best level falls below $-150$~dBm is subsequently overwritten with a
constant $-110$~dBm (\texttt{prepare\_\allowbreak{}transfer\_\allowbreak{}dataset\_\allowbreak{}v19.py:\allowbreak{}274}),
which introduces a point mass in the RSSI distribution that must be
excluded before any distributional claim about the target field is made.

\subsection{Coverage definition}

Coverage is a hard indicator rather than a probability, defined on the
serving-cell received level by a single threshold,
\begin{equation}
c_{i} = \mathbf{1}\bigl[\mathrm{RSSI}_{i} > -85\ \mathrm{dBm}\bigr] ,
\label{eq:app-coverage}
\end{equation}
% prepare_transfer_dataset_v19.py:275
so that $H(c_{i}\mid \mathrm{RSSI}_{i}) = 0$ exactly and the fifth
supervised channel carries no information beyond the fourth. There is no
fading margin, no outage probability, no shadowing standard deviation and
no receiver sensitivity or noise-figure calculation entering this
definition; the threshold is a fixed design constant of the generator. The
learned coverage channel is accordingly a derived classification of the
received level rather than an independent coverage probability.

\subsection{Antenna configuration}

Table~\ref{tab:app-antenna-params} separates the antenna quantities the
generator actually consumes from those that exist in the licensing extract
and are discarded, and Table~\ref{tab:app-antenna-inventory} gives the
resulting per-city inventory. The essential point for a reader assessing
the physical realism of the target is that the radiation pattern is
isotropic: azimuth, mechanical and electrical downtilt, half-power
beamwidth, front-to-back ratio and polarisation are all present in the
source records and none is read by the loader
(\texttt{prepare\_\allowbreak{}transfer\_\allowbreak{}dataset\_\allowbreak{}v19.py:\allowbreak{}73--79}). Because sectors of the
same site share coordinates, discarding azimuth turns each site into a
stack of co-located omnidirectional sources that differ only in frequency,
power, gain and height, with a mean of 13.5 and a maximum of 21 such
sources per site in Lins and a maximum of 49 in Campinas; the best-server
rule of Eq.~\eqref{eq:app-serving} then reduces each stack to whichever
member happens to maximise $P_{j}+G_{j}-L_{\mathrm{tot}}$. The loader also
retains catalogue gains of exactly 0~dBi where the licensing record carries
no value, affecting 32 transmitters in Campinas, 4 in Sorocaba and 1 in
Bauru, and it does not apply the 12~dBi substitution that the standalone
coverage script uses for the same situation
(\texttt{generate\_\allowbreak{}realistic\_\allowbreak{}coverage\allowbreak{}.py:\allowbreak{}191}), so the two entry points to
the same physics are not numerically identical.

\begin{table}[!t]
\caption{Antenna quantities in the target generator, grouped by whether
the model consumes them. Column names in typewriter font are fields of the
ANATEL licensing extract. The second group is the reason the radiation
pattern is isotropic: those fields are populated in the records and the
loader never reads them.}
\label{tab:app-antenna-params}
\centering
\footnotesize
\begin{tabular}{@{}ll@{}}
\toprule
Quantity & Source and unit in the model \\
\midrule
\multicolumn{2}{@{}l}{\textit{Used}} \\
Carrier frequency $f$   & \texttt{FreqTxMHz}, MHz \\
Transmit power $P$      & $10\log_{10}(1000P^{\mathrm{W}})$, dBm \\
Antenna gain $G$        & \texttt{GanhoAntena}, dBi, isotropic \\
Antenna height $h_{b}$  & \texttt{AlturaAntena}, m AGL \\
Site coordinates        & \texttt{Latitude}, \texttt{Longitude}, deg \\
Environment class $e$   & derived from $f$ alone \\
Mobile height $h_{m}$   & fixed at 1.5\,m for all nodes \\
\midrule
\multicolumn{2}{@{}l}{\textit{Present in the records but never read}} \\
Azimuth                 & \texttt{Azimute}, deg \\
Downtilt                & \texttt{AnguloElevacao}, deg \\
Half-power beamwidth    & \texttt{AnguloMeiaPotenciaAntena}, deg \\
Front-to-back ratio     & \texttt{FrenteCostaAntena}, dB \\
Polarisation            & \texttt{Polarizacao} \\
\midrule
\multicolumn{2}{@{}l}{\textit{Absent from the model altogether}} \\
Feeder, connector, body loss & --- \\
Receive antenna gain         & --- \\
Noise figure, thermal noise  & --- \\
Interference, SINR           & --- \\
Fading, shadowing margin     & --- \\
\bottomrule
\end{tabular}
\end{table}

\begin{table}[!t]
\caption{Per-city transmitter inventory after the generator's validity
filter ($f>0$ and $P^{\mathrm{W}}>0$). Sites are unique coordinate pairs.
The last column is the share of transmitters whose carrier exceeds the
2000\,MHz internal clip.}
\label{tab:app-antenna-inventory}
\centering
\footnotesize
\setlength{\tabcolsep}{4pt}
\begin{tabular}{@{}lrrccrr@{}}
\toprule
City & Tx & Sites & $f$ (MHz) & $P$ (dBm) & $h_{b}$ (m) & Clip \\
\midrule
Lins     &  175 &  13 & 788--2680 & 41.6--49.0 & 21--75 & 27.4\,\% \\
Campinas & 4252 & 350 & 778--3550 & 33.0--53.0 &  3--74 & 26.8\,\% \\
Sorocaba & 1398 & 120 & 788--2680 & 35.1--49.3 &  3--76 & 24.5\,\% \\
Bauru    &  794 &  75 & 788--2680 & 40.0--49.0 &  3--90 & 30.1\,\% \\
\bottomrule
\end{tabular}
\end{table}

\subsection{The two auxiliary channels}

The five supervised channels written by the V19 pipeline are total path
loss, two placeholder channels, serving-cell RSSI and the coverage
indicator; the placeholders are written as zeros at that stage
(\texttt{prepare\_\allowbreak{}transfer\_\allowbreak{}dataset\_\allowbreak{}v19.py:\allowbreak{}276--282}) and are then filled
by the enrichment step, which is part of the pipeline for every dataset
used in this study. That step computes a vegetation shadow margin
following ITU-R P.833 from the Sentinel-2 NDVI of the node and a
knife-edge diffraction loss following ITU-R P.526 from the DEM-derived
terrain ruggedness index and roughness
(\texttt{data\_\allowbreak{}raw/\allowbreak{}enrich\_\allowbreak{}rf\_\allowbreak{}targets\allowbreak{}.py:\allowbreak{}336--396,533--534}), using a
specific attenuation $\gamma = 0.2\,\mathrm{NDVI}\,(f_{\mathrm{GHz}}/1.8)^{0.4}$
in dB/m over a depth $d_{v} = d\,\operatorname{clip}(\mathrm{NDVI}-0.5,0,1)$
capped at 40~dB, and a single knife-edge obstacle placed at the path
midpoint with an obstruction height
$h_{\mathrm{obs}} = \operatorname{clip}(0.5\,\mathrm{TRI}_{\mathrm{m}} + 0.5\,r_{\mathrm{m}},\,0,\,80)$
in metres, also capped at 40~dB, with roughness $r_{\mathrm{m}}$ in metres after
correction. These two channels are therefore the point
at which the digital elevation model and the vegetation index enter the
reference field as physically labelled attenuation terms, and they are the
targets against which the decoder's vegetation and terrain heads are
trained and scored. Two properties of the enrichment must be stated because
they differ from the geometry of the total-loss channel: the distance it
uses is the distance to the nearest transmitter rather than to the serving
cell of Eq.~\eqref{eq:app-serving}
(\texttt{enrich\_\allowbreak{}rf\_\allowbreak{}targets\allowbreak{}.py:\allowbreak{}546}), and the frequency is a single
global 1800~MHz (\texttt{FREQ\_\allowbreak{}MHZ\_\allowbreak{}DEFAULT}, line 38) rather than the
per-transmitter carrier. The auxiliary channels are consequently
consistent with each other but evaluated under a different link geometry
and frequency from the total-loss channel, and the manuscript states this
explicitly where the three channels are compared.

\subsection{Target corrector}
\label{app:corrector}

Three defects of the enriched files were corrected before the spatially
blocked evaluation, all by \texttt{contrafactual\_\allowbreak{}alvo\_\allowbreak{}completo\allowbreak{}.py}, which
writes new files and leaves the originals untouched. First, the terrain term
of the total-loss channel used the random slope draw described above; the
corrector substitutes the node's DEM slope and updates the total-loss and
RSSI channels by exact difference, which is possible because
Eq.~\eqref{eq:app-aterrain} depends only on the node and does not change which
transmitter serves it. Second, the enrichment step read the topographic
position index column where the diffraction term expects the terrain
ruggedness index; the corrector reads the ruggedness column. Third, the
obstruction height added ruggedness in metres to roughness in normalised
units, whose scale differs by city; the corrector converts roughness to
metres with the per-city elevation standard deviation. The diffraction
channel is then recomputed with the original enrichment function. Sentinel
nodes keep their sentinel values in the total-loss and RSSI channels. The
16 corrected reference fields are released with their SHA-256 digests.

\subsection{What the generator does not model}

An emulation claim is only as strong as the honesty of the description of
what is being emulated, and several mechanisms that a radio-planning reader
would expect are simply absent here. There is no ray tracing and no terrain
profile along the link: the base term never sees an elevation sample
between transmitter and receiver, so diffraction over real obstacles,
Fresnel-zone clearance, line-of-sight testing, multipath and ground
reflection have no representation, and the knife-edge term of the
enrichment substitutes a single midpoint obstacle whose height is a
ruggedness proxy rather than a profile maximum. There is no clutter
database, as discussed above, and no building layer, so the urban,
suburban and rural corrections act as three global constants selected by
carrier frequency. There is no stochastic component of any kind in the
delivered path loss: the docstring of the propagation routine announces
log-normal fading (\texttt{generate\_\allowbreak{}realistic\_\allowbreak{}coverage\allowbreak{}.py:\allowbreak{}80}) but no
random term appears in the body of the function, and the only randomness in
the pipeline was the terrain slope draw described above, which the corrector
removes from the released fields. There is no noise floor, no
receiver noise figure, no interference from neighbouring cells and
therefore no SINR, so the received level of Eq.~\eqref{eq:app-rssi} is a
pure two-term link budget and coverage is a threshold on that budget rather
than a link-quality decision. Finally, the geometry is two-dimensional and
planar: node and transmitter altitudes never enter, the receiver is fixed
at 1.5~m above an unspecified reference surface, Earth curvature is
ignored, and the distance uses a constant metres-per-degree factor. Stating
these absences precisely is what allows the reformulated contribution to be
defended, because the quantity the network reproduces is a well-defined and
fully specified analytical field, and the accuracy figures in the body of
the paper measure emulation fidelity to that field and nothing more.
 immediately before
% the bibliography in main.tex. Requires amsmath and booktabs, both
% already loaded in the preamble (main.tex:11,16).
%
% PROVENANCE DISCIPLINE: every numerical value printed by this file is
% listed in the commented provenance table at the end, with the artefact
% field and the command that produced it. No number was read off a log.
% Artefact: REVISAO_R1/audit_target_generator.json
% Script:   REVISAO_R1/audit_target_generator.py
% =====================================================================

\section{Target Generation Model}
\label{app:target-generator}

\subsection{Software identification and scope}

The reference fields used throughout this study are not field measurements
but the output of a deterministic propagation generator written for this
project, and the contribution evaluated in this paper is
therefore the emulation of that generator at graph scale rather than the
prediction of measured radio conditions. No commercial or third-party
planning tool participates in the pipeline: there is no Atoll, ICS Telecom,
WinProp, TIREM or Longley--Rice component anywhere in the chain, and the
entire target field is produced by three Python modules and a post-processing corrector that depend, for
the propagation arithmetic itself, only on NumPy. Because those modules
carry no version string, they are identified here by the SHA-256 digest of
their contents, which is the only unambiguous version key available:
\texttt{generate\_\allowbreak{}realistic\_\allowbreak{}coverage\allowbreak{}.py} (\texttt{95ea0423\,2280\,6fd6},
13\,631~bytes) holds the propagation equations,
\texttt{01\_\allowbreak{}data/\allowbreak{}prepare\_\allowbreak{}transfer\_\allowbreak{}dataset\_\allowbreak{}v19.py}
(\texttt{45c16d43\,bc13\,a4bc}, 22\,917~bytes) applies them over the terrain
graph and writes the five supervised channels, and
\texttt{data\_\allowbreak{}raw/\allowbreak{}enrich\_\allowbreak{}rf\_\allowbreak{}targets\allowbreak{}.py} (\texttt{5d38012e\,c66e\,842d},
26\,012~bytes) replaces the two auxiliary channels with terrain-derived
attenuation terms. The corrector \texttt{contrafactual\_\allowbreak{}alvo\_\allowbreak{}completo\allowbreak{}.py}
(\texttt{ebeaf759\,2511\,fae3}, 10\,076~bytes) then rewrote the enriched files as
described in Section~\ref{app:corrector}; every result of the spatially blocked
evaluation uses the corrected files, whereas the comparisons quoted from the
earlier campaign use the enriched files before correction. Antenna
parameters come from a single ANATEL licensing extract,
\texttt{antenas\_\allowbreak{}interior\_\allowbreak{}sp\_\allowbreak{}final\allowbreak{}.csv}
(\texttt{ad9addb0\,c515\,a0b4}, 17\,763\,282~bytes). Everything reported in
this appendix is transcribed from those five files, and the equations below
are annotated with the source line at which each expression is implemented
so that the correspondence can be checked line by line.

\subsection{Base path loss}

The base term is a frequency-switched pair of empirical median-path-loss
models, both evaluated in decibels with distance in kilometres and
frequency in megahertz. The mobile-height correction factor for a
medium-to-small city is applied identically in both branches,
\begin{equation}
a(h_{m}) = \bigl(1.1\log_{10} f - 0.7\bigr)h_{m}
         - \bigl(1.56\log_{10} f - 0.8\bigr) ,
\label{eq:app-ahm}
\end{equation}
% generate_realistic_coverage.py:38 (Okumura branch) and :62 (COST-231 branch)
with the mobile height fixed at $h_{m}=1.5$~m for every node in the domain
(\texttt{prepare\_\allowbreak{}transfer\_\allowbreak{}dataset\_\allowbreak{}v19.py:\allowbreak{}210}). The Okumura--Hata urban
reference curve is
\begin{equation}
\begin{split}
L_{\mathrm{urb}}^{\mathrm{OH}} = {}& 69.55 + 26.16\log_{10} f
   - 13.82\log_{10} h_{b} - a(h_{m}) \\
 & + \bigl(44.9 - 6.55\log_{10} h_{b}\bigr)\log_{10} d ,
\end{split}
\label{eq:app-hata-urban}
\end{equation}
% generate_realistic_coverage.py:41-42
and the environment corrections applied to it are
\begin{equation}
L_{\mathrm{base}}^{\mathrm{OH}} =
\begin{cases}
L_{\mathrm{urb}}^{\mathrm{OH}}, & \text{urban},\\[2pt]
L_{\mathrm{urb}}^{\mathrm{OH}} - 2\bigl(\log_{10}\tfrac{f}{28}\bigr)^{2} - 5.4,
  & \text{suburban},\\[2pt]
L_{\mathrm{urb}}^{\mathrm{OH}} - 4.78\bigl(\log_{10} f\bigr)^{2} & \\
  \qquad {} + 18.33\log_{10} f - 40.94, & \text{rural}.
\end{cases}
\label{eq:app-hata-env}
\end{equation}

% generate_realistic_coverage.py:44-51
Above 1500~MHz the generator switches to the COST\,231 Hata extension,
\begin{equation}
\begin{split}
L_{\mathrm{base}}^{\mathrm{C231}} = {}& 46.3 + 33.9\log_{10} f
   - 13.82\log_{10} h_{b} - a(h_{m}) \\
 & + \bigl(44.9 - 6.55\log_{10} h_{b}\bigr)\log_{10} d + C_{m} ,
\end{split}
\label{eq:app-cost231}
\end{equation}
% generate_realistic_coverage.py:67-68; C_m at :65
where $C_{m}=3$~dB for the urban class and $C_{m}=0$~dB otherwise, and the
selection between the two branches is the hard threshold
$L_{\mathrm{base}} = L_{\mathrm{base}}^{\mathrm{OH}}$ if $f\le 1500$~MHz and
$L_{\mathrm{base}} = L_{\mathrm{base}}^{\mathrm{C231}}$ otherwise
(\texttt{generate\_\allowbreak{}realistic\_\allowbreak{}coverage\allowbreak{}.py:\allowbreak{}85--88}).

Three numerical guards inside these branches change the effective inputs
and must be declared because they are not part of either published model.
Frequency is clipped to $[150, 2000]$~MHz in the Okumura branch and to
$[1500, 2000]$~MHz in the COST\,231 branch
(\texttt{generate\_\allowbreak{}realistic\_\allowbreak{}coverage\allowbreak{}.py:\allowbreak{}34,58}), so that every carrier
above 2000~MHz enters the base term as if it were at exactly 2000~MHz,
while the vegetation term of Eq.~\eqref{eq:app-aveg} continues to receive
the true carrier (\texttt{generate\_\allowbreak{}realistic\_\allowbreak{}coverage\allowbreak{}.py:\allowbreak{}101}) and the two
components of the same sum are consequently evaluated at different
frequencies; in the four cities this clip is active for 27.4\,\% of the
Lins transmitters, 26.8\,\% of Campinas, 24.5\,\% of Sorocaba and 30.1\,\%
of Bauru, which means that for roughly a quarter of the field the base
frequency has been silently substituted. Distance is floored at $10^{-3}$~km,
that is at one metre (\texttt{generate\_\allowbreak{}realistic\_\allowbreak{}coverage\allowbreak{}.py:\allowbreak{}35,59}),
three decades below the 1~km lower bound of the Okumura--Hata validity
domain. Base-station height enters unclipped, and 5.1\,\% of the Lins
transmitters, 23.3\,\% of Campinas, 18.9\,\% of Sorocaba and 6.7\,\% of
Bauru sit below the 30~m floor of that same domain. The generator is
therefore an extrapolation of both empirical models over part of its input
range, and the reported accuracy figures should be read as fidelity to this
particular extrapolation rather than as agreement with the published
models inside their stated envelopes.

\subsection{Clutter and environment class}

What the generator calls the propagation environment is not a clutter
parameter in the usual sense, and this is the single point in the model
most likely to be misread. There is no land-cover raster, no morphology
database, no building or clutter-height layer, and no per-pixel clutter
correction anywhere in the chain; the class is assigned once per
transmitter by a threshold on its own carrier frequency, urban when
$f>1800$~MHz, suburban when $800<f\le 1800$~MHz and rural when
$f\le 800$~MHz (\texttt{prepare\_\allowbreak{}transfer\_\allowbreak{}dataset\_\allowbreak{}v19.py:\allowbreak{}84--86}). The
consequence is that the environment label is a deterministic function of
frequency and carries no independent geographic information, so that
Eq.~\eqref{eq:app-hata-env} does not describe the ground under a given node
but only the band of the transmitter serving it. A verifiable corollary of
this coupling is that the urban class can never reach the Okumura branch
and the rural class can never reach COST\,231, which the transmitter records
confirm exactly: zero transmitters fall in either impossible combination in any of
the four cities. Terrain clutter is therefore unmodelled, and a frequency-indexed
correction constant stands in its place.

\subsection{Vegetation attenuation}

Vegetation enters through a single additive term parameterised by NDVI
taken raw from column 12 of the terrain feature matrix
(\texttt{prepare\_\allowbreak{}transfer\_\allowbreak{}dataset\_\allowbreak{}v19.py:\allowbreak{}140}), which the code maps to a
density and then to an effective path depth,
\begin{equation}
\rho_{v} = \operatorname{clip}\!\left(\frac{\mathrm{NDVI}+1}{2},\,0,\,1\right),
\qquad
d_{v} = 0.3\, \rho_{v}\, d_{m} ,
\label{eq:app-vegdepth}
\end{equation}
% generate_realistic_coverage.py:95,98 -- d_m in metres
so that at most 30\,\% of the geometric path is declared to lie inside
vegetation regardless of where the vegetation actually is. The attenuation
follows the ITU-R P.833 functional form for a single vegetation entry,
\begin{equation}
A_{\mathrm{veg}} = \operatorname{clip}\!\bigl(
0.25 f^{0.39} \min(d_{v}, 400)^{0.25},\, 0,\, 40 \bigr)\ [\mathrm{dB}] ,
\label{eq:app-aveg}
\end{equation}
% generate_realistic_coverage.py:101-102 -- f in MHz, d_v in metres
with frequency in megahertz and depth in metres, a depth cap of 400~m and
an attenuation cap of 40~dB. The depth cap is the dominant feature of this
term in practice: for a node with $\mathrm{NDVI}=0.5$ the cap is reached at
a geometric distance of about 1778~m, beyond which $A_{\mathrm{veg}}$ is
constant in distance and depends only on frequency, taking the values
15.87~dB at 900~MHz, 20.80~dB at 1800~MHz and 24.00~dB at 2600~MHz.
Because the argument of the quarter-power is a product of distance and
NDVI, the term is also strongly distance-driven below the cap, reaching
7.73~dB at 100~m and 13.74~dB at 1~km for the same NDVI at 900~MHz, so that
most of its dynamic range is geometry rather than vegetation. The 40~dB
clip is unreachable in this dataset: at the highest carrier present in the
four cities, 3550~MHz, and with the depth cap and NDVI both saturated, the
term attains only 27.10~dB.

\subsection{Terrain attenuation}

The terrain term is the component of the generator that most needs an
explicit and uncomfortable statement, because its name does not describe
its content. Formally it is a piecewise-linear function of a dimensionless
slope variable,
\begin{equation}
A_{\mathrm{ter}} = \operatorname{clip}\!\Bigl(
20\,s + 30\,\max(s-0.1,\,0),\; 0,\; 30 \Bigr)\ [\mathrm{dB}] ,
\label{eq:app-aterrain}
\end{equation}
% generate_realistic_coverage.py:109-116
with $s$ the local slope expressed as a rise-over-run ratio. The V19 stage
of the pipeline did not read $s$ from the digital elevation model: it drew
it independently for each node from a uniform distribution on $[0.02, 0.15]$
by an unseeded call to the NumPy generator
(\texttt{prepare\_\allowbreak{}transfer\_\allowbreak{}dataset\_\allowbreak{}v19.py:\allowbreak{}141}, and identically at
\texttt{generate\_\allowbreak{}realistic\_\allowbreak{}coverage\allowbreak{}.py:\allowbreak{}167}), which made
Eq.~\eqref{eq:app-aterrain} an independent random offset per node between
0.4~dB and 4.5~dB, with mean 1.99~dB and standard deviation 1.15~dB over the
support of that distribution. The corrector of Section~\ref{app:corrector}
replaces the draw with the node's DEM slope, stored as feature~1 of the
released files, and recomputes the total-loss and RSSI channels by exact
difference, so that in every reference field used for the blocked evaluation
Eq.~\eqref{eq:app-aterrain} is a deterministic function of terrain.
Recomputing the corrected total-loss channel from the features stored in the
released files reproduces it with a mean absolute error of at most
$3.02 \times 10^{-4}$~dB per cell over 150\,000 covered nodes in five cells.

\subsection{Total loss, link budget and serving cell}

The three components are summed without interaction terms,
\begin{equation}
\begin{split}
L_{\mathrm{tot}}(i,j) = {}& L_{\mathrm{base}}(d_{ij}, f_{j}, h_{b,j}, e_{j}) \\
 & + A_{\mathrm{veg}}(d_{ij}, \mathrm{NDVI}_{i})
   + A_{\mathrm{ter}}(s_{i}) ,
\end{split}
\label{eq:app-ltot}
\end{equation}
% generate_realistic_coverage.py:121
and converted to a received level by a two-term link budget that contains
no cable, connector, body or feeder loss and no receive antenna gain,
\begin{equation}
\mathrm{RSSI}(i,j) = P_{j} + G_{j} - L_{\mathrm{tot}}(i,j)
\quad [\mathrm{dBm}] ,
\label{eq:app-rssi}
\end{equation}
% prepare_transfer_dataset_v19.py:215
where $P_{j}=10\log_{10}(1000\,P_{j}^{\mathrm{W}})$ is the transmitter power
in dBm (\texttt{prepare\_\allowbreak{}transfer\_\allowbreak{}dataset\_\allowbreak{}v19.py:\allowbreak{}81}) and $G_{j}$ is the
catalogue antenna gain in dBi, applied as a scalar in every direction. The
serving cell is then the strongest received level among the transmitters
whose planar distance to the node is within the 30~km inclusion radius,
\begin{equation}
s(i) = \!\!\operatorname*{arg\,max}_{j \,:\, d_{ij} \le 30\,\mathrm{km}}\!\!
\mathrm{RSSI}(i,j),
\quad
L_{i} = L_{\mathrm{tot}}\bigl(i, s(i)\bigr) ,
\label{eq:app-serving}
\end{equation}
% prepare_transfer_dataset_v19.py:220-228; radius at :46
which is a pure best-server rule with no handover hysteresis, no load
balancing, no cell-selection offset and no minimum-received-level
qualification. Distances are planar rather than geodesic, computed as
$\Delta y = 111000\,\Delta\varphi$ and
$\Delta x = 111000\,\Delta\lambda\cos\varphi$ in metres
(\texttt{prepare\_\allowbreak{}transfer\_\allowbreak{}dataset\_\allowbreak{}v19.py:\allowbreak{}183--185}); the constant
111000~m per degree understates the WGS-84 mean meridian degree of
111194.93~m by 0.175\,\%, a bias that is negligible against the decibel
scale of the model but is nonetheless a systematic contraction of every
distance in the dataset. Nodes with no transmitter inside the radius retain
the sentinel initialisation of $-200$~dBm and 300~dB, and any node whose
best level falls below $-150$~dBm is subsequently overwritten with a
constant $-110$~dBm (\texttt{prepare\_\allowbreak{}transfer\_\allowbreak{}dataset\_\allowbreak{}v19.py:\allowbreak{}274}),
which introduces a point mass in the RSSI distribution that must be
excluded before any distributional claim about the target field is made.

\subsection{Coverage definition}

Coverage is a hard indicator rather than a probability, defined on the
serving-cell received level by a single threshold,
\begin{equation}
c_{i} = \mathbf{1}\bigl[\mathrm{RSSI}_{i} > -85\ \mathrm{dBm}\bigr] ,
\label{eq:app-coverage}
\end{equation}
% prepare_transfer_dataset_v19.py:275
so that $H(c_{i}\mid \mathrm{RSSI}_{i}) = 0$ exactly and the fifth
supervised channel carries no information beyond the fourth. There is no
fading margin, no outage probability, no shadowing standard deviation and
no receiver sensitivity or noise-figure calculation entering this
definition; the threshold is a fixed design constant of the generator. The
learned coverage channel is accordingly a derived classification of the
received level rather than an independent coverage probability.

\subsection{Antenna configuration}

Table~\ref{tab:app-antenna-params} separates the antenna quantities the
generator actually consumes from those that exist in the licensing extract
and are discarded, and Table~\ref{tab:app-antenna-inventory} gives the
resulting per-city inventory. The essential point for a reader assessing
the physical realism of the target is that the radiation pattern is
isotropic: azimuth, mechanical and electrical downtilt, half-power
beamwidth, front-to-back ratio and polarisation are all present in the
source records and none is read by the loader
(\texttt{prepare\_\allowbreak{}transfer\_\allowbreak{}dataset\_\allowbreak{}v19.py:\allowbreak{}73--79}). Because sectors of the
same site share coordinates, discarding azimuth turns each site into a
stack of co-located omnidirectional sources that differ only in frequency,
power, gain and height, with a mean of 13.5 and a maximum of 21 such
sources per site in Lins and a maximum of 49 in Campinas; the best-server
rule of Eq.~\eqref{eq:app-serving} then reduces each stack to whichever
member happens to maximise $P_{j}+G_{j}-L_{\mathrm{tot}}$. The loader also
retains catalogue gains of exactly 0~dBi where the licensing record carries
no value, affecting 32 transmitters in Campinas, 4 in Sorocaba and 1 in
Bauru, and it does not apply the 12~dBi substitution that the standalone
coverage script uses for the same situation
(\texttt{generate\_\allowbreak{}realistic\_\allowbreak{}coverage\allowbreak{}.py:\allowbreak{}191}), so the two entry points to
the same physics are not numerically identical.

\begin{table}[!t]
\caption{Antenna quantities in the target generator, grouped by whether
the model consumes them. Column names in typewriter font are fields of the
ANATEL licensing extract. The second group is the reason the radiation
pattern is isotropic: those fields are populated in the records and the
loader never reads them.}
\label{tab:app-antenna-params}
\centering
\footnotesize
\begin{tabular}{@{}ll@{}}
\toprule
Quantity & Source and unit in the model \\
\midrule
\multicolumn{2}{@{}l}{\textit{Used}} \\
Carrier frequency $f$   & \texttt{FreqTxMHz}, MHz \\
Transmit power $P$      & $10\log_{10}(1000P^{\mathrm{W}})$, dBm \\
Antenna gain $G$        & \texttt{GanhoAntena}, dBi, isotropic \\
Antenna height $h_{b}$  & \texttt{AlturaAntena}, m AGL \\
Site coordinates        & \texttt{Latitude}, \texttt{Longitude}, deg \\
Environment class $e$   & derived from $f$ alone \\
Mobile height $h_{m}$   & fixed at 1.5\,m for all nodes \\
\midrule
\multicolumn{2}{@{}l}{\textit{Present in the records but never read}} \\
Azimuth                 & \texttt{Azimute}, deg \\
Downtilt                & \texttt{AnguloElevacao}, deg \\
Half-power beamwidth    & \texttt{AnguloMeiaPotenciaAntena}, deg \\
Front-to-back ratio     & \texttt{FrenteCostaAntena}, dB \\
Polarisation            & \texttt{Polarizacao} \\
\midrule
\multicolumn{2}{@{}l}{\textit{Absent from the model altogether}} \\
Feeder, connector, body loss & --- \\
Receive antenna gain         & --- \\
Noise figure, thermal noise  & --- \\
Interference, SINR           & --- \\
Fading, shadowing margin     & --- \\
\bottomrule
\end{tabular}
\end{table}

\begin{table}[!t]
\caption{Per-city transmitter inventory after the generator's validity
filter ($f>0$ and $P^{\mathrm{W}}>0$). Sites are unique coordinate pairs.
The last column is the share of transmitters whose carrier exceeds the
2000\,MHz internal clip.}
\label{tab:app-antenna-inventory}
\centering
\footnotesize
\setlength{\tabcolsep}{4pt}
\begin{tabular}{@{}lrrccrr@{}}
\toprule
City & Tx & Sites & $f$ (MHz) & $P$ (dBm) & $h_{b}$ (m) & Clip \\
\midrule
Lins     &  175 &  13 & 788--2680 & 41.6--49.0 & 21--75 & 27.4\,\% \\
Campinas & 4252 & 350 & 778--3550 & 33.0--53.0 &  3--74 & 26.8\,\% \\
Sorocaba & 1398 & 120 & 788--2680 & 35.1--49.3 &  3--76 & 24.5\,\% \\
Bauru    &  794 &  75 & 788--2680 & 40.0--49.0 &  3--90 & 30.1\,\% \\
\bottomrule
\end{tabular}
\end{table}

\subsection{The two auxiliary channels}

The five supervised channels written by the V19 pipeline are total path
loss, two placeholder channels, serving-cell RSSI and the coverage
indicator; the placeholders are written as zeros at that stage
(\texttt{prepare\_\allowbreak{}transfer\_\allowbreak{}dataset\_\allowbreak{}v19.py:\allowbreak{}276--282}) and are then filled
by the enrichment step, which is part of the pipeline for every dataset
used in this study. That step computes a vegetation shadow margin
following ITU-R P.833 from the Sentinel-2 NDVI of the node and a
knife-edge diffraction loss following ITU-R P.526 from the DEM-derived
terrain ruggedness index and roughness
(\texttt{data\_\allowbreak{}raw/\allowbreak{}enrich\_\allowbreak{}rf\_\allowbreak{}targets\allowbreak{}.py:\allowbreak{}336--396,533--534}), using a
specific attenuation $\gamma = 0.2\,\mathrm{NDVI}\,(f_{\mathrm{GHz}}/1.8)^{0.4}$
in dB/m over a depth $d_{v} = d\,\operatorname{clip}(\mathrm{NDVI}-0.5,0,1)$
capped at 40~dB, and a single knife-edge obstacle placed at the path
midpoint with an obstruction height
$h_{\mathrm{obs}} = \operatorname{clip}(0.5\,\mathrm{TRI}_{\mathrm{m}} + 0.5\,r_{\mathrm{m}},\,0,\,80)$
in metres, also capped at 40~dB, with roughness $r_{\mathrm{m}}$ in metres after
correction. These two channels are therefore the point
at which the digital elevation model and the vegetation index enter the
reference field as physically labelled attenuation terms, and they are the
targets against which the decoder's vegetation and terrain heads are
trained and scored. Two properties of the enrichment must be stated because
they differ from the geometry of the total-loss channel: the distance it
uses is the distance to the nearest transmitter rather than to the serving
cell of Eq.~\eqref{eq:app-serving}
(\texttt{enrich\_\allowbreak{}rf\_\allowbreak{}targets\allowbreak{}.py:\allowbreak{}546}), and the frequency is a single
global 1800~MHz (\texttt{FREQ\_\allowbreak{}MHZ\_\allowbreak{}DEFAULT}, line 38) rather than the
per-transmitter carrier. The auxiliary channels are consequently
consistent with each other but evaluated under a different link geometry
and frequency from the total-loss channel, and the manuscript states this
explicitly where the three channels are compared.

\subsection{Target corrector}
\label{app:corrector}

Three defects of the enriched files were corrected before the spatially
blocked evaluation, all by \texttt{contrafactual\_\allowbreak{}alvo\_\allowbreak{}completo\allowbreak{}.py}, which
writes new files and leaves the originals untouched. First, the terrain term
of the total-loss channel used the random slope draw described above; the
corrector substitutes the node's DEM slope and updates the total-loss and
RSSI channels by exact difference, which is possible because
Eq.~\eqref{eq:app-aterrain} depends only on the node and does not change which
transmitter serves it. Second, the enrichment step read the topographic
position index column where the diffraction term expects the terrain
ruggedness index; the corrector reads the ruggedness column. Third, the
obstruction height added ruggedness in metres to roughness in normalised
units, whose scale differs by city; the corrector converts roughness to
metres with the per-city elevation standard deviation. The diffraction
channel is then recomputed with the original enrichment function. Sentinel
nodes keep their sentinel values in the total-loss and RSSI channels. The
16 corrected reference fields are released with their SHA-256 digests.

\subsection{What the generator does not model}

An emulation claim is only as strong as the honesty of the description of
what is being emulated, and several mechanisms that a radio-planning reader
would expect are simply absent here. There is no ray tracing and no terrain
profile along the link: the base term never sees an elevation sample
between transmitter and receiver, so diffraction over real obstacles,
Fresnel-zone clearance, line-of-sight testing, multipath and ground
reflection have no representation, and the knife-edge term of the
enrichment substitutes a single midpoint obstacle whose height is a
ruggedness proxy rather than a profile maximum. There is no clutter
database, as discussed above, and no building layer, so the urban,
suburban and rural corrections act as three global constants selected by
carrier frequency. There is no stochastic component of any kind in the
delivered path loss: the docstring of the propagation routine announces
log-normal fading (\texttt{generate\_\allowbreak{}realistic\_\allowbreak{}coverage\allowbreak{}.py:\allowbreak{}80}) but no
random term appears in the body of the function, and the only randomness in
the pipeline was the terrain slope draw described above, which the corrector
removes from the released fields. There is no noise floor, no
receiver noise figure, no interference from neighbouring cells and
therefore no SINR, so the received level of Eq.~\eqref{eq:app-rssi} is a
pure two-term link budget and coverage is a threshold on that budget rather
than a link-quality decision. Finally, the geometry is two-dimensional and
planar: node and transmitter altitudes never enter, the receiver is fixed
at 1.5~m above an unspecified reference surface, Earth curvature is
ignored, and the distance uses a constant metres-per-degree factor. Stating
these absences precisely is what allows the reformulated contribution to be
defended, because the quantity the network reproduces is a well-defined and
fully specified analytical field, and the accuracy figures in the body of
the paper measure emulation fidelity to that field and nothing more.
 immediately before
% the bibliography in main.tex. Requires amsmath and booktabs, both
% already loaded in the preamble (main.tex:11,16).
%
% PROVENANCE DISCIPLINE: every numerical value printed by this file is
% listed in the commented provenance table at the end, with the artefact
% field and the command that produced it. No number was read off a log.
% Artefact: REVISAO_R1/audit_target_generator.json
% Script:   REVISAO_R1/audit_target_generator.py
% =====================================================================

\section{Target Generation Model}
\label{app:target-generator}

\subsection{Software identification and scope}

The reference fields used throughout this study are not field measurements
but the output of a deterministic propagation generator written for this
project, and the contribution evaluated in this paper is
therefore the emulation of that generator at graph scale rather than the
prediction of measured radio conditions. No commercial or third-party
planning tool participates in the pipeline: there is no Atoll, ICS Telecom,
WinProp, TIREM or Longley--Rice component anywhere in the chain, and the
entire target field is produced by three Python modules and a post-processing corrector that depend, for
the propagation arithmetic itself, only on NumPy. Because those modules
carry no version string, they are identified here by the SHA-256 digest of
their contents, which is the only unambiguous version key available:
\texttt{generate\_\allowbreak{}realistic\_\allowbreak{}coverage\allowbreak{}.py} (\texttt{95ea0423\,2280\,6fd6},
13\,631~bytes) holds the propagation equations,
\texttt{01\_\allowbreak{}data/\allowbreak{}prepare\_\allowbreak{}transfer\_\allowbreak{}dataset\_\allowbreak{}v19.py}
(\texttt{45c16d43\,bc13\,a4bc}, 22\,917~bytes) applies them over the terrain
graph and writes the five supervised channels, and
\texttt{data\_\allowbreak{}raw/\allowbreak{}enrich\_\allowbreak{}rf\_\allowbreak{}targets\allowbreak{}.py} (\texttt{5d38012e\,c66e\,842d},
26\,012~bytes) replaces the two auxiliary channels with terrain-derived
attenuation terms. The corrector \texttt{contrafactual\_\allowbreak{}alvo\_\allowbreak{}completo\allowbreak{}.py}
(\texttt{ebeaf759\,2511\,fae3}, 10\,076~bytes) then rewrote the enriched files as
described in Section~\ref{app:corrector}; every result of the spatially blocked
evaluation uses the corrected files, whereas the comparisons quoted from the
earlier campaign use the enriched files before correction. Antenna
parameters come from a single ANATEL licensing extract,
\texttt{antenas\_\allowbreak{}interior\_\allowbreak{}sp\_\allowbreak{}final\allowbreak{}.csv}
(\texttt{ad9addb0\,c515\,a0b4}, 17\,763\,282~bytes). Everything reported in
this appendix is transcribed from those five files, and the equations below
are annotated with the source line at which each expression is implemented
so that the correspondence can be checked line by line.

\subsection{Base path loss}

The base term is a frequency-switched pair of empirical median-path-loss
models, both evaluated in decibels with distance in kilometres and
frequency in megahertz. The mobile-height correction factor for a
medium-to-small city is applied identically in both branches,
\begin{equation}
a(h_{m}) = \bigl(1.1\log_{10} f - 0.7\bigr)h_{m}
         - \bigl(1.56\log_{10} f - 0.8\bigr) ,
\label{eq:app-ahm}
\end{equation}
% generate_realistic_coverage.py:38 (Okumura branch) and :62 (COST-231 branch)
with the mobile height fixed at $h_{m}=1.5$~m for every node in the domain
(\texttt{prepare\_\allowbreak{}transfer\_\allowbreak{}dataset\_\allowbreak{}v19.py:\allowbreak{}210}). The Okumura--Hata urban
reference curve is
\begin{equation}
\begin{split}
L_{\mathrm{urb}}^{\mathrm{OH}} = {}& 69.55 + 26.16\log_{10} f
   - 13.82\log_{10} h_{b} - a(h_{m}) \\
 & + \bigl(44.9 - 6.55\log_{10} h_{b}\bigr)\log_{10} d ,
\end{split}
\label{eq:app-hata-urban}
\end{equation}
% generate_realistic_coverage.py:41-42
and the environment corrections applied to it are
\begin{equation}
L_{\mathrm{base}}^{\mathrm{OH}} =
\begin{cases}
L_{\mathrm{urb}}^{\mathrm{OH}}, & \text{urban},\\[2pt]
L_{\mathrm{urb}}^{\mathrm{OH}} - 2\bigl(\log_{10}\tfrac{f}{28}\bigr)^{2} - 5.4,
  & \text{suburban},\\[2pt]
L_{\mathrm{urb}}^{\mathrm{OH}} - 4.78\bigl(\log_{10} f\bigr)^{2} & \\
  \qquad {} + 18.33\log_{10} f - 40.94, & \text{rural}.
\end{cases}
\label{eq:app-hata-env}
\end{equation}

% generate_realistic_coverage.py:44-51
Above 1500~MHz the generator switches to the COST\,231 Hata extension,
\begin{equation}
\begin{split}
L_{\mathrm{base}}^{\mathrm{C231}} = {}& 46.3 + 33.9\log_{10} f
   - 13.82\log_{10} h_{b} - a(h_{m}) \\
 & + \bigl(44.9 - 6.55\log_{10} h_{b}\bigr)\log_{10} d + C_{m} ,
\end{split}
\label{eq:app-cost231}
\end{equation}
% generate_realistic_coverage.py:67-68; C_m at :65
where $C_{m}=3$~dB for the urban class and $C_{m}=0$~dB otherwise, and the
selection between the two branches is the hard threshold
$L_{\mathrm{base}} = L_{\mathrm{base}}^{\mathrm{OH}}$ if $f\le 1500$~MHz and
$L_{\mathrm{base}} = L_{\mathrm{base}}^{\mathrm{C231}}$ otherwise
(\texttt{generate\_\allowbreak{}realistic\_\allowbreak{}coverage\allowbreak{}.py:\allowbreak{}85--88}).

Three numerical guards inside these branches change the effective inputs
and must be declared because they are not part of either published model.
Frequency is clipped to $[150, 2000]$~MHz in the Okumura branch and to
$[1500, 2000]$~MHz in the COST\,231 branch
(\texttt{generate\_\allowbreak{}realistic\_\allowbreak{}coverage\allowbreak{}.py:\allowbreak{}34,58}), so that every carrier
above 2000~MHz enters the base term as if it were at exactly 2000~MHz,
while the vegetation term of Eq.~\eqref{eq:app-aveg} continues to receive
the true carrier (\texttt{generate\_\allowbreak{}realistic\_\allowbreak{}coverage\allowbreak{}.py:\allowbreak{}101}) and the two
components of the same sum are consequently evaluated at different
frequencies; in the four cities this clip is active for 27.4\,\% of the
Lins transmitters, 26.8\,\% of Campinas, 24.5\,\% of Sorocaba and 30.1\,\%
of Bauru, which means that for roughly a quarter of the field the base
frequency has been silently substituted. Distance is floored at $10^{-3}$~km,
that is at one metre (\texttt{generate\_\allowbreak{}realistic\_\allowbreak{}coverage\allowbreak{}.py:\allowbreak{}35,59}),
three decades below the 1~km lower bound of the Okumura--Hata validity
domain. Base-station height enters unclipped, and 5.1\,\% of the Lins
transmitters, 23.3\,\% of Campinas, 18.9\,\% of Sorocaba and 6.7\,\% of
Bauru sit below the 30~m floor of that same domain. The generator is
therefore an extrapolation of both empirical models over part of its input
range, and the reported accuracy figures should be read as fidelity to this
particular extrapolation rather than as agreement with the published
models inside their stated envelopes.

\subsection{Clutter and environment class}

What the generator calls the propagation environment is not a clutter
parameter in the usual sense, and this is the single point in the model
most likely to be misread. There is no land-cover raster, no morphology
database, no building or clutter-height layer, and no per-pixel clutter
correction anywhere in the chain; the class is assigned once per
transmitter by a threshold on its own carrier frequency, urban when
$f>1800$~MHz, suburban when $800<f\le 1800$~MHz and rural when
$f\le 800$~MHz (\texttt{prepare\_\allowbreak{}transfer\_\allowbreak{}dataset\_\allowbreak{}v19.py:\allowbreak{}84--86}). The
consequence is that the environment label is a deterministic function of
frequency and carries no independent geographic information, so that
Eq.~\eqref{eq:app-hata-env} does not describe the ground under a given node
but only the band of the transmitter serving it. A verifiable corollary of
this coupling is that the urban class can never reach the Okumura branch
and the rural class can never reach COST\,231, which the transmitter records
confirm exactly: zero transmitters fall in either impossible combination in any of
the four cities. Terrain clutter is therefore unmodelled, and a frequency-indexed
correction constant stands in its place.

\subsection{Vegetation attenuation}

Vegetation enters through a single additive term parameterised by NDVI
taken raw from column 12 of the terrain feature matrix
(\texttt{prepare\_\allowbreak{}transfer\_\allowbreak{}dataset\_\allowbreak{}v19.py:\allowbreak{}140}), which the code maps to a
density and then to an effective path depth,
\begin{equation}
\rho_{v} = \operatorname{clip}\!\left(\frac{\mathrm{NDVI}+1}{2},\,0,\,1\right),
\qquad
d_{v} = 0.3\, \rho_{v}\, d_{m} ,
\label{eq:app-vegdepth}
\end{equation}
% generate_realistic_coverage.py:95,98 -- d_m in metres
so that at most 30\,\% of the geometric path is declared to lie inside
vegetation regardless of where the vegetation actually is. The attenuation
follows the ITU-R P.833 functional form for a single vegetation entry,
\begin{equation}
A_{\mathrm{veg}} = \operatorname{clip}\!\bigl(
0.25 f^{0.39} \min(d_{v}, 400)^{0.25},\, 0,\, 40 \bigr)\ [\mathrm{dB}] ,
\label{eq:app-aveg}
\end{equation}
% generate_realistic_coverage.py:101-102 -- f in MHz, d_v in metres
with frequency in megahertz and depth in metres, a depth cap of 400~m and
an attenuation cap of 40~dB. The depth cap is the dominant feature of this
term in practice: for a node with $\mathrm{NDVI}=0.5$ the cap is reached at
a geometric distance of about 1778~m, beyond which $A_{\mathrm{veg}}$ is
constant in distance and depends only on frequency, taking the values
15.87~dB at 900~MHz, 20.80~dB at 1800~MHz and 24.00~dB at 2600~MHz.
Because the argument of the quarter-power is a product of distance and
NDVI, the term is also strongly distance-driven below the cap, reaching
7.73~dB at 100~m and 13.74~dB at 1~km for the same NDVI at 900~MHz, so that
most of its dynamic range is geometry rather than vegetation. The 40~dB
clip is unreachable in this dataset: at the highest carrier present in the
four cities, 3550~MHz, and with the depth cap and NDVI both saturated, the
term attains only 27.10~dB.

\subsection{Terrain attenuation}

The terrain term is the component of the generator that most needs an
explicit and uncomfortable statement, because its name does not describe
its content. Formally it is a piecewise-linear function of a dimensionless
slope variable,
\begin{equation}
A_{\mathrm{ter}} = \operatorname{clip}\!\Bigl(
20\,s + 30\,\max(s-0.1,\,0),\; 0,\; 30 \Bigr)\ [\mathrm{dB}] ,
\label{eq:app-aterrain}
\end{equation}
% generate_realistic_coverage.py:109-116
with $s$ the local slope expressed as a rise-over-run ratio. The V19 stage
of the pipeline did not read $s$ from the digital elevation model: it drew
it independently for each node from a uniform distribution on $[0.02, 0.15]$
by an unseeded call to the NumPy generator
(\texttt{prepare\_\allowbreak{}transfer\_\allowbreak{}dataset\_\allowbreak{}v19.py:\allowbreak{}141}, and identically at
\texttt{generate\_\allowbreak{}realistic\_\allowbreak{}coverage\allowbreak{}.py:\allowbreak{}167}), which made
Eq.~\eqref{eq:app-aterrain} an independent random offset per node between
0.4~dB and 4.5~dB, with mean 1.99~dB and standard deviation 1.15~dB over the
support of that distribution. The corrector of Section~\ref{app:corrector}
replaces the draw with the node's DEM slope, stored as feature~1 of the
released files, and recomputes the total-loss and RSSI channels by exact
difference, so that in every reference field used for the blocked evaluation
Eq.~\eqref{eq:app-aterrain} is a deterministic function of terrain.
Recomputing the corrected total-loss channel from the features stored in the
released files reproduces it with a mean absolute error of at most
$3.02 \times 10^{-4}$~dB per cell over 150\,000 covered nodes in five cells.

\subsection{Total loss, link budget and serving cell}

The three components are summed without interaction terms,
\begin{equation}
\begin{split}
L_{\mathrm{tot}}(i,j) = {}& L_{\mathrm{base}}(d_{ij}, f_{j}, h_{b,j}, e_{j}) \\
 & + A_{\mathrm{veg}}(d_{ij}, \mathrm{NDVI}_{i})
   + A_{\mathrm{ter}}(s_{i}) ,
\end{split}
\label{eq:app-ltot}
\end{equation}
% generate_realistic_coverage.py:121
and converted to a received level by a two-term link budget that contains
no cable, connector, body or feeder loss and no receive antenna gain,
\begin{equation}
\mathrm{RSSI}(i,j) = P_{j} + G_{j} - L_{\mathrm{tot}}(i,j)
\quad [\mathrm{dBm}] ,
\label{eq:app-rssi}
\end{equation}
% prepare_transfer_dataset_v19.py:215
where $P_{j}=10\log_{10}(1000\,P_{j}^{\mathrm{W}})$ is the transmitter power
in dBm (\texttt{prepare\_\allowbreak{}transfer\_\allowbreak{}dataset\_\allowbreak{}v19.py:\allowbreak{}81}) and $G_{j}$ is the
catalogue antenna gain in dBi, applied as a scalar in every direction. The
serving cell is then the strongest received level among the transmitters
whose planar distance to the node is within the 30~km inclusion radius,
\begin{equation}
s(i) = \!\!\operatorname*{arg\,max}_{j \,:\, d_{ij} \le 30\,\mathrm{km}}\!\!
\mathrm{RSSI}(i,j),
\quad
L_{i} = L_{\mathrm{tot}}\bigl(i, s(i)\bigr) ,
\label{eq:app-serving}
\end{equation}
% prepare_transfer_dataset_v19.py:220-228; radius at :46
which is a pure best-server rule with no handover hysteresis, no load
balancing, no cell-selection offset and no minimum-received-level
qualification. Distances are planar rather than geodesic, computed as
$\Delta y = 111000\,\Delta\varphi$ and
$\Delta x = 111000\,\Delta\lambda\cos\varphi$ in metres
(\texttt{prepare\_\allowbreak{}transfer\_\allowbreak{}dataset\_\allowbreak{}v19.py:\allowbreak{}183--185}); the constant
111000~m per degree understates the WGS-84 mean meridian degree of
111194.93~m by 0.175\,\%, a bias that is negligible against the decibel
scale of the model but is nonetheless a systematic contraction of every
distance in the dataset. Nodes with no transmitter inside the radius retain
the sentinel initialisation of $-200$~dBm and 300~dB, and any node whose
best level falls below $-150$~dBm is subsequently overwritten with a
constant $-110$~dBm (\texttt{prepare\_\allowbreak{}transfer\_\allowbreak{}dataset\_\allowbreak{}v19.py:\allowbreak{}274}),
which introduces a point mass in the RSSI distribution that must be
excluded before any distributional claim about the target field is made.

\subsection{Coverage definition}

Coverage is a hard indicator rather than a probability, defined on the
serving-cell received level by a single threshold,
\begin{equation}
c_{i} = \mathbf{1}\bigl[\mathrm{RSSI}_{i} > -85\ \mathrm{dBm}\bigr] ,
\label{eq:app-coverage}
\end{equation}
% prepare_transfer_dataset_v19.py:275
so that $H(c_{i}\mid \mathrm{RSSI}_{i}) = 0$ exactly and the fifth
supervised channel carries no information beyond the fourth. There is no
fading margin, no outage probability, no shadowing standard deviation and
no receiver sensitivity or noise-figure calculation entering this
definition; the threshold is a fixed design constant of the generator. The
learned coverage channel is accordingly a derived classification of the
received level rather than an independent coverage probability.

\subsection{Antenna configuration}

Table~\ref{tab:app-antenna-params} separates the antenna quantities the
generator actually consumes from those that exist in the licensing extract
and are discarded, and Table~\ref{tab:app-antenna-inventory} gives the
resulting per-city inventory. The essential point for a reader assessing
the physical realism of the target is that the radiation pattern is
isotropic: azimuth, mechanical and electrical downtilt, half-power
beamwidth, front-to-back ratio and polarisation are all present in the
source records and none is read by the loader
(\texttt{prepare\_\allowbreak{}transfer\_\allowbreak{}dataset\_\allowbreak{}v19.py:\allowbreak{}73--79}). Because sectors of the
same site share coordinates, discarding azimuth turns each site into a
stack of co-located omnidirectional sources that differ only in frequency,
power, gain and height, with a mean of 13.5 and a maximum of 21 such
sources per site in Lins and a maximum of 49 in Campinas; the best-server
rule of Eq.~\eqref{eq:app-serving} then reduces each stack to whichever
member happens to maximise $P_{j}+G_{j}-L_{\mathrm{tot}}$. The loader also
retains catalogue gains of exactly 0~dBi where the licensing record carries
no value, affecting 32 transmitters in Campinas, 4 in Sorocaba and 1 in
Bauru, and it does not apply the 12~dBi substitution that the standalone
coverage script uses for the same situation
(\texttt{generate\_\allowbreak{}realistic\_\allowbreak{}coverage\allowbreak{}.py:\allowbreak{}191}), so the two entry points to
the same physics are not numerically identical.

\begin{table}[!t]
\caption{Antenna quantities in the target generator, grouped by whether
the model consumes them. Column names in typewriter font are fields of the
ANATEL licensing extract. The second group is the reason the radiation
pattern is isotropic: those fields are populated in the records and the
loader never reads them.}
\label{tab:app-antenna-params}
\centering
\footnotesize
\begin{tabular}{@{}ll@{}}
\toprule
Quantity & Source and unit in the model \\
\midrule
\multicolumn{2}{@{}l}{\textit{Used}} \\
Carrier frequency $f$   & \texttt{FreqTxMHz}, MHz \\
Transmit power $P$      & $10\log_{10}(1000P^{\mathrm{W}})$, dBm \\
Antenna gain $G$        & \texttt{GanhoAntena}, dBi, isotropic \\
Antenna height $h_{b}$  & \texttt{AlturaAntena}, m AGL \\
Site coordinates        & \texttt{Latitude}, \texttt{Longitude}, deg \\
Environment class $e$   & derived from $f$ alone \\
Mobile height $h_{m}$   & fixed at 1.5\,m for all nodes \\
\midrule
\multicolumn{2}{@{}l}{\textit{Present in the records but never read}} \\
Azimuth                 & \texttt{Azimute}, deg \\
Downtilt                & \texttt{AnguloElevacao}, deg \\
Half-power beamwidth    & \texttt{AnguloMeiaPotenciaAntena}, deg \\
Front-to-back ratio     & \texttt{FrenteCostaAntena}, dB \\
Polarisation            & \texttt{Polarizacao} \\
\midrule
\multicolumn{2}{@{}l}{\textit{Absent from the model altogether}} \\
Feeder, connector, body loss & --- \\
Receive antenna gain         & --- \\
Noise figure, thermal noise  & --- \\
Interference, SINR           & --- \\
Fading, shadowing margin     & --- \\
\bottomrule
\end{tabular}
\end{table}

\begin{table}[!t]
\caption{Per-city transmitter inventory after the generator's validity
filter ($f>0$ and $P^{\mathrm{W}}>0$). Sites are unique coordinate pairs.
The last column is the share of transmitters whose carrier exceeds the
2000\,MHz internal clip.}
\label{tab:app-antenna-inventory}
\centering
\footnotesize
\setlength{\tabcolsep}{4pt}
\begin{tabular}{@{}lrrccrr@{}}
\toprule
City & Tx & Sites & $f$ (MHz) & $P$ (dBm) & $h_{b}$ (m) & Clip \\
\midrule
Lins     &  175 &  13 & 788--2680 & 41.6--49.0 & 21--75 & 27.4\,\% \\
Campinas & 4252 & 350 & 778--3550 & 33.0--53.0 &  3--74 & 26.8\,\% \\
Sorocaba & 1398 & 120 & 788--2680 & 35.1--49.3 &  3--76 & 24.5\,\% \\
Bauru    &  794 &  75 & 788--2680 & 40.0--49.0 &  3--90 & 30.1\,\% \\
\bottomrule
\end{tabular}
\end{table}

\subsection{The two auxiliary channels}

The five supervised channels written by the V19 pipeline are total path
loss, two placeholder channels, serving-cell RSSI and the coverage
indicator; the placeholders are written as zeros at that stage
(\texttt{prepare\_\allowbreak{}transfer\_\allowbreak{}dataset\_\allowbreak{}v19.py:\allowbreak{}276--282}) and are then filled
by the enrichment step, which is part of the pipeline for every dataset
used in this study. That step computes a vegetation shadow margin
following ITU-R P.833 from the Sentinel-2 NDVI of the node and a
knife-edge diffraction loss following ITU-R P.526 from the DEM-derived
terrain ruggedness index and roughness
(\texttt{data\_\allowbreak{}raw/\allowbreak{}enrich\_\allowbreak{}rf\_\allowbreak{}targets\allowbreak{}.py:\allowbreak{}336--396,533--534}), using a
specific attenuation $\gamma = 0.2\,\mathrm{NDVI}\,(f_{\mathrm{GHz}}/1.8)^{0.4}$
in dB/m over a depth $d_{v} = d\,\operatorname{clip}(\mathrm{NDVI}-0.5,0,1)$
capped at 40~dB, and a single knife-edge obstacle placed at the path
midpoint with an obstruction height
$h_{\mathrm{obs}} = \operatorname{clip}(0.5\,\mathrm{TRI}_{\mathrm{m}} + 0.5\,r_{\mathrm{m}},\,0,\,80)$
in metres, also capped at 40~dB, with roughness $r_{\mathrm{m}}$ in metres after
correction. These two channels are therefore the point
at which the digital elevation model and the vegetation index enter the
reference field as physically labelled attenuation terms, and they are the
targets against which the decoder's vegetation and terrain heads are
trained and scored. Two properties of the enrichment must be stated because
they differ from the geometry of the total-loss channel: the distance it
uses is the distance to the nearest transmitter rather than to the serving
cell of Eq.~\eqref{eq:app-serving}
(\texttt{enrich\_\allowbreak{}rf\_\allowbreak{}targets\allowbreak{}.py:\allowbreak{}546}), and the frequency is a single
global 1800~MHz (\texttt{FREQ\_\allowbreak{}MHZ\_\allowbreak{}DEFAULT}, line 38) rather than the
per-transmitter carrier. The auxiliary channels are consequently
consistent with each other but evaluated under a different link geometry
and frequency from the total-loss channel, and the manuscript states this
explicitly where the three channels are compared.

\subsection{Target corrector}
\label{app:corrector}

Three defects of the enriched files were corrected before the spatially
blocked evaluation, all by \texttt{contrafactual\_\allowbreak{}alvo\_\allowbreak{}completo\allowbreak{}.py}, which
writes new files and leaves the originals untouched. First, the terrain term
of the total-loss channel used the random slope draw described above; the
corrector substitutes the node's DEM slope and updates the total-loss and
RSSI channels by exact difference, which is possible because
Eq.~\eqref{eq:app-aterrain} depends only on the node and does not change which
transmitter serves it. Second, the enrichment step read the topographic
position index column where the diffraction term expects the terrain
ruggedness index; the corrector reads the ruggedness column. Third, the
obstruction height added ruggedness in metres to roughness in normalised
units, whose scale differs by city; the corrector converts roughness to
metres with the per-city elevation standard deviation. The diffraction
channel is then recomputed with the original enrichment function. Sentinel
nodes keep their sentinel values in the total-loss and RSSI channels. The
16 corrected reference fields are released with their SHA-256 digests.

\subsection{What the generator does not model}

An emulation claim is only as strong as the honesty of the description of
what is being emulated, and several mechanisms that a radio-planning reader
would expect are simply absent here. There is no ray tracing and no terrain
profile along the link: the base term never sees an elevation sample
between transmitter and receiver, so diffraction over real obstacles,
Fresnel-zone clearance, line-of-sight testing, multipath and ground
reflection have no representation, and the knife-edge term of the
enrichment substitutes a single midpoint obstacle whose height is a
ruggedness proxy rather than a profile maximum. There is no clutter
database, as discussed above, and no building layer, so the urban,
suburban and rural corrections act as three global constants selected by
carrier frequency. There is no stochastic component of any kind in the
delivered path loss: the docstring of the propagation routine announces
log-normal fading (\texttt{generate\_\allowbreak{}realistic\_\allowbreak{}coverage\allowbreak{}.py:\allowbreak{}80}) but no
random term appears in the body of the function, and the only randomness in
the pipeline was the terrain slope draw described above, which the corrector
removes from the released fields. There is no noise floor, no
receiver noise figure, no interference from neighbouring cells and
therefore no SINR, so the received level of Eq.~\eqref{eq:app-rssi} is a
pure two-term link budget and coverage is a threshold on that budget rather
than a link-quality decision. Finally, the geometry is two-dimensional and
planar: node and transmitter altitudes never enter, the receiver is fixed
at 1.5~m above an unspecified reference surface, Earth curvature is
ignored, and the distance uses a constant metres-per-degree factor. Stating
these absences precisely is what allows the reformulated contribution to be
defended, because the quantity the network reproduces is a well-defined and
fully specified analytical field, and the accuracy figures in the body of
the paper measure emulation fidelity to that field and nothing more.
